% =====================================================================
%  ThaiGer Track Control - Printable Trilingual Guidebook
%  -------------------------------------------------------------------
%  ENGINE: XeLaTeX  (required for Chinese / CJK and system fonts)
%
%  Build:
%      xelatex guidebook.tex
%      xelatex guidebook.tex        (run twice for the table of contents)
%  or:
%      latexmk -xelatex guidebook.tex
%
%  FONTS: The Chinese part uses Windows fonts (SimSun / Microsoft YaHei).
%  On Linux/macOS replace them with Noto (see the \setCJK... lines below).
% =====================================================================
\documentclass[11pt,a4paper,oneside]{report}

% --- Fonts & engine ---
\usepackage{fontspec}
% Main text font must cover Latin + Cyrillic (for the Russian part AND its TOC /
% bookmark entries, which are typeset outside the Russian language switch).
% Times New Roman ships with Windows; on Linux/macOS use "TeX Gyre Termes" or
% "Noto Serif" (both bundled with TeX / widely available and Cyrillic-capable).
\setmainfont{Times New Roman}
\usepackage{xeCJK}               % Chinese support (load before babel-driven text use)

% Windows 11 CJK fonts. Linux/macOS alternative shown in comments:
%   \setCJKmainfont{Noto Serif CJK SC}
%   \setCJKsansfont{Noto Sans CJK SC}
%   \setCJKmonofont{Noto Sans Mono CJK SC}
\setCJKmainfont{SimSun}
\setCJKsansfont{Microsoft YaHei}
\setCJKmonofont{Microsoft YaHei}

% --- Languages (hyphenation) ---
% shorthands=off so babel never makes ", :, <, >, ~ active - they appear in
% code tokens / URLs and must stay literal.
\usepackage[main=english,ngerman,french,spanish,russian,shorthands=off]{babel}

% --- Layout ---
\usepackage[a4paper,margin=2.4cm,headheight=15pt]{geometry}
\usepackage{parskip}
\usepackage{textcomp}
\usepackage{pifont}
\usepackage{amssymb}
\usepackage{xcolor}
\usepackage{array}
\usepackage{booktabs}
\usepackage{tabularx}
\usepackage{enumitem}
\usepackage{titlesec}
\usepackage{fancyhdr}
\usepackage{listings}
\usepackage{tcolorbox}
\usepackage{hyperref}
\usepackage{bookmark}            % cleaner PDF bookmarks (part -> section jump)

% Let TeX absorb small over-runs instead of pushing into the margin
\setlength{\emergencystretch}{3.5em}

% --- Brand colours (print-friendly) ---
\definecolor{brandblue}{HTML}{2B5CE6}
\definecolor{ink}{HTML}{1A1F26}
\definecolor{tgreen}{HTML}{0F9D58}
\definecolor{tamber}{HTML}{C97A00}
\definecolor{tred}{HTML}{D7263D}
\definecolor{codebg}{HTML}{F4F6F9}
\definecolor{codeframe}{HTML}{E1E7EF}

\hypersetup{
  colorlinks=true, linkcolor=brandblue, urlcolor=brandblue, citecolor=brandblue,
  pdftitle={ThaiGer Track Control - Guidebook},
  pdfauthor={ThaiGer Racing},
  pdfsubject={Telemetry cockpit user guide},
}

% --- Section numbering: per-part 1..N (no chapter prefix) ---
\renewcommand{\thesection}{\arabic{section}}
\titleformat{\section}{\Large\bfseries\color{brandblue}}{\thesection.}{0.5em}{}
\titleformat{\subsection}{\large\bfseries\color{ink}}{\thesubsection}{0.5em}{}
\titleformat{\subsubsection}{\normalsize\bfseries\color{ink}}{}{0em}{}
\titlespacing*{\section}{0pt}{14pt}{6pt}

% --- Header / footer ---
\pagestyle{fancy}
\fancyhf{}
\fancyhead[L]{\small\textcolor{ink}{ThaiGer Track Control}}
\fancyhead[R]{\small\textcolor{ink}{User Guidebook}}
\fancyfoot[C]{\small\thepage}
\renewcommand{\headrulewidth}{0.4pt}
\renewcommand{\footrulewidth}{0pt}

% --- Lists ---
\setlist{itemsep=2pt,topsep=3pt,parsep=0pt,leftmargin=1.5em}

% --- Code listings ---
\lstdefinestyle{tg}{
  backgroundcolor=\color{codebg},
  basicstyle=\ttfamily\small,
  breaklines=true,
  frame=single,
  rulecolor=\color{codeframe},
  framesep=6pt,
  xleftmargin=8pt, xrightmargin=8pt,
  columns=fullflexible,
  keepspaces=true,
  showstringspaces=false,
  aboveskip=8pt, belowskip=8pt,
}
\lstset{style=tg}

% --- Callout box ---
\newtcolorbox{tgnote}{colback=blue!4,colframe=brandblue,boxrule=0.5pt,arc=2pt,
  left=8pt,right=8pt,top=6pt,bottom=6pt}

% --- Wrapping table column ---
\newcolumntype{L}{>{\raggedright\arraybackslash}X}

% --- Symbol helpers (avoid bare Unicode in Latin text) ---
\newcommand{\arrow}{\,$\rightarrow$\,}
\newcommand{\degC}{\textdegree C}
\newcommand{\Htwo}{H\textsubscript{2}}
\newcommand{\midsep}{\,\textperiodcentered\,}
\newcommand{\ok}{\textcolor{tgreen}{\ding{51}}}
\newcommand{\nope}{\textcolor{tred}{\ding{55}}}
\newcommand{\code}[1]{\texttt{#1}}

% Part heading helper
\newcommand{\guidepart}[1]{%
  \clearpage
  \part*{#1}%
  \addcontentsline{toc}{part}{#1}%
  \setcounter{section}{0}%
}

\begin{document}

% =====================================================================
%  TITLE PAGE
% =====================================================================
\begin{titlepage}
\centering
\vspace*{2.5cm}
{\Huge\bfseries\color{brandblue} ThaiGer \Htwo{} Track Control\par}
\vspace{0.6cm}
{\LARGE User Guidebook\par}
\vspace{0.3cm}
{\large Telemetry cockpit for hydrogen fuel-cell racing\par}
\vspace{1.4cm}
{\large Thaiger 7 \midsep{} Bengalo\par}
\vspace{2.2cm}
\begin{tcolorbox}[width=0.7\linewidth,colback=codebg,colframe=codeframe,
  boxrule=0.5pt,arc=3pt,halign=center]
\large
English \midsep{} Deutsch \midsep{} Français\\
Español \midsep{} Русский \midsep{} 中文 (简体)
\end{tcolorbox}
\vfill
{\small Android cockpit app \,+\, browser-based engineer dashboard\par}
{\small Document generated for print \midsep{} A4\par}
\end{titlepage}

% =====================================================================
%  TABLE OF CONTENTS
% =====================================================================
\clearpage
\tableofcontents
\clearpage


% #####################################################################
%  PART I - ENGLISH
% #####################################################################
\selectlanguage{english}
\guidepart{English Guide}

\section{What this project is}
ThaiGer Track Control is a two-part telemetry system for a hydrogen fuel-cell
race car:
\begin{itemize}
  \item An \textbf{Android app} that runs on a phone \emph{inside the car}. It
    receives live telemetry from the car's controller over \textbf{Bluetooth Low
    Energy (BLE)}, shows it on a race dashboard for the driver, records the run,
    and \textbf{relays} the data over the internet via \textbf{MQTT}.
  \item An \textbf{Engineer Dashboard} --- a single HTML file
    (\code{engineer\_dashboard.html}) --- that runs in a browser on the
    \emph{engineer's laptop} in the pit. It subscribes to the same MQTT stream and
    shows operations data, a live map with a power-colored GPS trail, and
    analytics charts.
\end{itemize}
The two halves never talk directly; they meet at the \textbf{MQTT broker} in the
cloud.

\section{System overview}

\begin{lstlisting}
   +------------+   BLE    +------------------+   MQTT/TLS   +----------------+
   | Car        | -------> | Android app      | -----------> | MQTT broker    |
   | controller | *A28.3*  | (phone in car)   | JSON frames  | (HiveMQ Cloud) |
   +------------+          +------------------+              +-------+--------+
                                  |                                  | WebSocket/TLS
                                  | GPS (phone LocationManager)      v
                                  v                          +----------------+
                          thaiger/{car}/gps                  | Engineer       |
                          thaiger/{car}/telemetry            | Dashboard HTML |
                                                             +----------------+
\end{lstlisting}
\begin{itemize}
  \item Two MQTT topics per car: \code{thaiger/\{car\}/telemetry} and
    \code{thaiger/\{car\}/gps}.
  \item \code{\{car\}} is the vehicle id: \code{thaiger7} or \code{bengalo}.
  \item The phone publishes; the engineer dashboard subscribes. Any number of
    laptops can watch the same car at once.
\end{itemize}

\section{The Android app}

\subsection{Installation \& first start}
The app is a standard Android project (Gradle). To build a debug APK:

\begin{lstlisting}
./gradlew assembleDebug
\end{lstlisting}

\begin{tabularx}{\linewidth}{@{}lL@{}}
\toprule
\textbf{Item} & \textbf{Value}\\
\midrule
Minimum Android & 9.0 (API 28)\\
Target Android  & 14 (API 34)\\
Language        & Java 11\\
App id          & \code{com.thaiger.h2racing}\\
\bottomrule
\end{tabularx}

\medskip
Install the APK on the phone that lives in the car. On first launch the app opens
on the \textbf{Car Select} screen.

\begin{tgnote}
\textbf{Tip.} Mount the phone in landscape where the driver can see it. The
dashboard keeps the screen awake during a run (Settings \arrow{} \emph{Keep screen
on}).
\end{tgnote}

\subsection{Permissions}
The app asks for permissions on the \textbf{Connecting} screen, before any
hardware is used. Grant all of them for full functionality.

\begin{tabularx}{\linewidth}{@{}lL@{}}
\toprule
\textbf{Permission} & \textbf{Why it is needed}\\
\midrule
Bluetooth Connect / Scan (Android 12+) & Pair \& talk to the car's BLE module\\
Bluetooth / Bluetooth Admin (Android $\le$ 11) & Same, on older Android\\
Location (fine) & Required by Android for BLE scanning \emph{and} for GPS\\
Location (coarse) & GPS fallback\\
Internet & MQTT relay to the engineer laptop\\
Wake Lock & Keep the screen on during a run\\
Vibrate & Haptic feedback for alerts\\
\bottomrule
\end{tabularx}

\medskip
If you deny Bluetooth, you can still use \textbf{Demo Mode} (synthetic data) to
explore the app.

\subsection{Screen 1 --- Car Select}
\begin{itemize}
  \item Two cards: \textbf{Thaiger 7} and \textbf{Bengalo}. Tap a card to select
    it (a blue outline marks the selection). Thaiger 7 is the default.
  \item The selected profile controls warning thresholds and the max-power scale.
  \item Tap \textbf{Settings} (top) to open app settings before a run.
  \item Tap \textbf{Connect to \ldots} to continue to the Connecting screen.
\end{itemize}

\subsection{Screen 2 --- Connecting (Bluetooth / Demo Mode)}
This screen pairs the phone with the car and starts the data stream.
\begin{enumerate}
  \item The app checks permissions and that Bluetooth is on.
  \item A dialog lists your \textbf{paired Bluetooth devices} plus a final
    \textbf{Demo Mode} entry.
  \item Choose your car's module (e.g. an HM-10 / HC-05 / ESP32-named device),
    \textbf{or} pick \textbf{Demo Mode} to run with a synthetic data stream and no
    hardware.
  \item \textbf{Pair New} opens the Android Bluetooth settings if your module is
    not paired yet. \textbf{Cancel} returns to Car Select.
\end{enumerate}
When the module reports \textbf{Connected}, the app automatically advances to the
Live Dashboard. If the MQTT relay is enabled in Settings, it starts here too, so
the engineer laptop begins receiving data as soon as the run starts.

\begin{tgnote}
Connection drops are handled automatically with exponential backoff
(1 \arrow{} 2 \arrow{} 4 \arrow{} 8 \arrow{} 16 s). You will see \emph{Reconnecting\ldots}
--- no action needed.
\end{tgnote}

\subsection{Screen 3 --- Live Dashboard}
The dashboard is the driver's main view (landscape, screen stays awake). It shows
the latest value of every telemetry field, refreshed continuously.
\begin{itemize}
  \item \textbf{Speed} --- the large hero number (km/h). If \emph{Speed color
    warning} is on, the number turns \textbf{red} when you drop below the car's
    minimum target speed and stays light otherwise. A small indicator shows
    \ok{} on target / \nope{} too slow.
  \item \textbf{Lap overlay} (single line above the bottom bar):
    \code{PREV 1:18} \midsep{} \code{LAP 3} \midsep{} \code{CURR 0:09} --- previous
    lap time, current lap number, and the running current-lap time.
  \item \textbf{Electrical block} --- fuel-cell / supercap / motor voltage and
    current, own consumption, cell voltage difference.
  \item \textbf{Thermal \& efficiency} --- FC temperature, air-pump duty, FC and
    system efficiency, energies, optimal speed, driving hint.
  \item \textbf{Alert banner} --- appears when a threshold is crossed (e.g. FC
    temperature or cell voltage difference too high). With \emph{Vibrate} /
    \emph{Alert sound} on you also get haptic / audio feedback. Tap
    \textbf{Dismiss} to hide the banner.
\end{itemize}
\begin{tgnote}
\textbf{Ending a run:} \textbf{long-press the speed number.} This stops Bluetooth
and the MQTT relay, finalizes the run statistics, and opens the \textbf{Post-Run
Summary}.
\end{tgnote}

\subsection{Screen 4 --- Post-Run Summary \& CSV export}
After a run you see aggregated results: run duration, distance, average speed,
energy used, peak motor power, max FC temperature, and alert count, plus a
\textbf{power graph} for the whole run.
\begin{itemize}
  \item \textbf{Export CSV} --- writes every recorded telemetry frame to a CSV
    file and opens the Android share sheet (email, Drive, messenger, \ldots). The
    file is named \code{thaiger\_\{car\}\_\allowbreak\{YYYYMMDD\_HHMM\}.csv} and contains 21
    columns. If no data was recorded you get a \emph{``No data to export''}
    message.
  \item \textbf{New Run} --- go straight back to Connecting for another run.
  \item \textbf{Back to Car Select} --- return to the start.
\end{itemize}

\subsection{Screen 5 --- Settings}
Reached from the \textbf{Settings} link on Car Select. All values are saved
immediately.

\subsubsection{Display}

\begin{tabularx}{\linewidth}{@{}llL@{}}
\toprule
\textbf{Setting} & \textbf{Default} & \textbf{Meaning}\\
\midrule
Keep screen on & On & Hold a wake lock so the screen never sleeps during a run\\
Speed color warning & On & Turn the speed red when below the car's minimum target speed\\
Update rate & 50 ms & How often the dashboard UI redraws (20--500 ms)\\
\bottomrule
\end{tabularx}

\subsubsection{Alerts}

\begin{tabularx}{\linewidth}{@{}llL@{}}
\toprule
\textbf{Setting} & \textbf{Default} & \textbf{Meaning}\\
\midrule
Vibrate & On & Haptic pulse when an alert fires\\
Alert sound & Off & Play a tone when an alert fires\\
\bottomrule
\end{tabularx}

\subsubsection{Thresholds (per vehicle)}
Stored separately for Thaiger 7 and Bengalo:

\begin{tabularx}{\linewidth}{@{}llL@{}}
\toprule
\textbf{Setting} & \textbf{Range} & \textbf{Default (Thaiger 7 / Bengalo)}\\
\midrule
FC temperature max & 30--120 \degC{} & 70 \degC{} / 65 \degC{}\\
Cell voltage diff max & 10--500 mV & 50 mV / 50 mV\\
\bottomrule
\end{tabularx}

\subsection{Screen 6 --- MQTT Settings}
A sub-screen of Settings that configures the relay to the engineer laptop. The
relay only runs if it is enabled \emph{before} you start a run.

\begin{tabularx}{\linewidth}{@{}llL@{}}
\toprule
\textbf{Field} & \textbf{Default} & \textbf{Notes}\\
\midrule
Enable relay & Off & Master switch --- turn on to publish telemetry\\
Use TLS & On & Encrypted connection (required by HiveMQ Cloud)\\
Broker host & --- & e.g. \code{abc123.s1.eu.hivemq.cloud}\\
Port & 8883 & TLS MQTT port (8883 typical; 1883 plain)\\
Username & --- & Broker username\\
Password & --- & Broker password (shown masked)\\
Publish rate & 5 Hz & Telemetry publish rate, 1--10 Hz\\
\bottomrule
\end{tabularx}
\begin{tgnote}
The phone publishes over \textbf{native MQTT (TLS)}. The engineer dashboard
connects to the \emph{same} broker over \textbf{MQTT-over-WebSocket} --- so the
broker must also expose a WebSocket listener (HiveMQ Cloud does, on port 8884).
\end{tgnote}

\section{The Engineer Dashboard}
\code{engineer\_dashboard.html} is a single self-contained file. It needs no build
step and no server --- just a modern browser and internet access (for the map
tiles, fonts, and the MQTT/charts libraries from CDN).

\subsection{Opening \& connecting}
\begin{enumerate}
  \item Open \code{engineer\_dashboard.html} in Chrome/Edge/Firefox.
  \item The \textbf{Connection Settings} dialog opens automatically on first use
    (or click the gear button, top-right).
  \item Fill in: \textbf{Car} (\code{Thaiger 7} or \code{Bengalo}, must match the
    car the phone publishes as); \textbf{Broker URL (WebSocket)}, e.g.
    \code{wss://abc123.s1.eu.hivemq.cloud:\allowbreak8884/mqtt} (use \code{wss://} for TLS, or
    \code{ws://host:port} for a plain local broker); \textbf{Username} /
    \textbf{Password}.
  \item Click \textbf{Connect}. Settings are saved in the browser
    (\code{localStorage}) and reconnect automatically next time.
\end{enumerate}
The status dot (top-right) shows the link: grey disconnected, amber connecting,
green connected. Three tabs switch between \textbf{Operations},
\textbf{Live Tracking}, and \textbf{Analytics}.

\subsection{View 1 --- Operations}
A single-glance overview: a \textbf{hero speed card} (blue at/above optimal speed,
red below) with average speed, lap count, run time, target lap and distance; an
\textbf{Electrical} grid; a \textbf{Thermal \& Efficiency} grid; and a pulsing
\textbf{LIVE} pill confirming fresh data.

\subsection{View 2 --- Live Tracking \& map}
A full-screen dark \textbf{OpenStreetMap} (Leaflet) following the car.
\begin{itemize}
  \item \textbf{Power-colored trail} --- the path is drawn segment by segment and
    colored by the motor power at that moment: \textcolor{tgreen}{green (0 W)}
    \arrow{} \textcolor{tamber}{amber (100 W)} \arrow{} \textcolor{tred}{red
    (200 W+)}. The newest part is fully opaque and fades with age, so you read both
    \emph{where} the car went and \emph{how hard it was working}. (200 W is the
    ThaiGer max tracking power; the constant \code{POWER\_MAX} can be changed.)
  \item \textbf{Power legend} (bottom-left) --- the color scale plus a live watt
    readout.
  \item \textbf{Rotating car marker} points in the direction of travel.
  \item \textbf{GPS status panel} (right) --- fix quality, accuracy, latitude,
    longitude, altitude, GPS speed, heading, live motor power, lap count, and a
    \textbf{Filtered} counter (rejected noisy fixes --- see below).
\end{itemize}

\subsubsection{GPS jitter cleanup}
Raw GPS can jump around. Each fix is cleaned before it touches the map by four
steps: (1) drop fixes with accuracy worse than 75 m; (2) reject ``teleports'' ---
fixes implying an impossible speed (over approx.\ 162 km/h), with a re-sync guard
so it never gets stuck; (3) ignore sub-1.5 m wobble while stationary; (4) light
smoothing of the remaining noise. The result is a clean trail with no random
jumps.

\subsubsection{Per-lap snapshot + reset}
Every time the car's lap counter increments, the dashboard automatically renders
the just-finished lap's trail to a PNG image and downloads it
(\code{thaiger-lapN-<timestamp>.png}), then clears the trail so the new lap is
drawn fresh. A green toast confirms each save.

\subsection{View 3 --- Analytics}
A \textbf{live telemetry table} listing all 22 fields with their current values,
plus a \textbf{Speed} chart (km/h over time) and an \textbf{Efficiency} chart (FC
and system efficiency over time), each showing the last 60 samples.

\section{Telemetry protocol \& field reference}
The car controller sends ASCII frames over BLE. Each field is wrapped in asterisks
with a single-letter key:

\begin{lstlisting}
*A28.3**B25.0**C3**D15:35**N53.4*
\end{lstlisting}
The parser is stream-tolerant: it handles partial packets, multi-chunk frames and
missing fields, always keeping the latest known value of each field.

\begin{tabularx}{\linewidth}{@{}clL@{}}
\toprule
\textbf{Key} & \textbf{Field} & \textbf{Unit}\\
\midrule
A & Speed & km/h\\
B & Average speed & km/h\\
C & Lap count & ---\\
D & Total run time & mm:ss\\
E & Target lap time & mm:ss\\
F & Optimal speed & km/h\\
G & Fuel-cell voltage & V\\
H & Supercap voltage & V\\
I & Motor voltage & V\\
J & Fuel-cell current & A\\
K & Supercap current & A\\
L & Motor current & A\\
M & Own consumption & A\\
N & Fuel-cell temperature & \degC{}\\
O & Air-pump duty cycle & \%\\
P & Driving hint & ---\\
Q & Cell voltage difference & mV\\
R & Fuel-cell energy (cumulative) & Ws\\
S & Motor energy (cumulative) & Ws\\
T & Fuel-cell efficiency & \%\\
U & System efficiency & \%\\
\bottomrule
\end{tabularx}

\medskip
Distance, motor power ($V \times A$), and energy in Wh are derived on the phone.

\section{MQTT topics \& payloads}
All messages are \textbf{QoS 0} (fire-and-forget). Empty/NaN fields are omitted to
keep payloads small (about 100--200 bytes).

\noindent\textbf{\code{thaiger/\{car\}/telemetry}} --- published at the configured
rate (default 5 Hz):

\begin{lstlisting}
{"ts":1736517201234,"A":28.3,"B":25.0,"C":3,"D":"15:35","G":28.3,"N":53.4,"dist":0.42}
\end{lstlisting}
\noindent\textbf{\code{thaiger/\{car\}/gps}} --- published when a GPS fix is
available:

\begin{lstlisting}
{"ts":1736517201300,"lat":13.756331,"lon":100.501762,"alt":12.0,"spd":53.4,"acc":3.5,"hdg":270.0}
\end{lstlisting}

\begin{tabularx}{\linewidth}{@{}lLl@{}}
\toprule
\textbf{Field} & \textbf{Meaning} & \textbf{Unit}\\
\midrule
ts & Unix timestamp & ms\\
lat / lon & Latitude / longitude & deg\\
alt & Altitude (if available) & m\\
spd & GPS speed (if available) & km/h\\
acc & Horizontal accuracy (if available) & m\\
hdg & Heading / bearing (if available) & deg\\
\bottomrule
\end{tabularx}

\section{Vehicle profiles}

\begin{tabularx}{\linewidth}{@{}llllL@{}}
\toprule
\textbf{Profile} & \textbf{id} & \textbf{Max power} & \textbf{FC temp limit} & \textbf{Min target speed}\\
\midrule
Thaiger 7 & \code{thaiger7} & 600 W & 70 \degC{} & 25 km/h\\
Bengalo   & \code{bengalo}  & 900 W & 65 \degC{} & 25 km/h\\
\bottomrule
\end{tabularx}

\medskip
Thresholds can be overridden per car in \textbf{Settings}. The engineer-dashboard
map trail uses a separate \textbf{200 W} power scale (\code{POWER\_MAX}), tuned for
the tracking display rather than the car's absolute peak.

\section{Troubleshooting}

\begin{tabularx}{\linewidth}{@{}LL@{}}
\toprule
\textbf{Symptom} & \textbf{Cause / fix}\\
\midrule
App won't connect over Bluetooth & Pair the module in Android Bluetooth settings first; check it is powered; grant Bluetooth + Location permissions. Use \textbf{Pair New}.\\
No data, but ``Connected'' & Wrong device, or controller not sending. Try Demo Mode to confirm the app, then re-check the module.\\
Dashboard never goes green & Wrong WebSocket URL (must be \code{wss://\ldots:8884/mqtt} for HiveMQ Cloud, not the \code{8883} native port), wrong credentials, or no WebSocket listener.\\
Dashboard green but no telemetry & The \textbf{Car} in the dashboard does not match the phone's car id, or the relay is not enabled on the phone. Both must use the same \code{thaiger7}/\code{bengalo}.\\
Map is blank & No internet for OSM tiles, or no GPS fixes yet. The map centers only once a fix arrives.\\
GPS trail jumps around & Cleanup should prevent this; if your device only has very coarse fixes, raise \code{GPS\_MAX\_ACC}. The \textbf{Filtered} counter shows how many fixes are dropped.\\
GPS updates only every few seconds & Ensure a clear sky view and that Location permission is granted; indoors it falls back to coarse network fixes.\\
No CSV produced & Complete a run with recorded frames; \emph{Export CSV} reports ``No data to export'' if nothing was captured.\\
\bottomrule
\end{tabularx}

\section{Tuning \& advanced settings}
On the phone (Settings / MQTT Settings): screen wake, speed color, UI update rate,
vibrate/sound, per-car thresholds, and all broker/relay parameters.

\medskip
\noindent In \code{engineer\_dashboard.html} (constants near the top of the script
block):

\begin{tabularx}{\linewidth}{@{}llL@{}}
\toprule
\textbf{Constant} & \textbf{Default} & \textbf{Purpose}\\
\midrule
\code{POWER\_MAX} & 200 & Watt value mapped to full red on the trail/legend\\
\code{GPS\_MAX\_ACC} & 75 & Drop GPS fixes worse than this many metres\\
\code{GPS\_MAX\_SPEED} & 45 & m/s; faster implied speed is treated as a teleport\\
\code{GPS\_MIN\_MOVE} & 1.5 & m; smaller moves are treated as stationary jitter\\
\code{GPS\_SMOOTH} & 0.35 & EMA smoothing factor (0 = none, 1 = raw)\\
\code{GPS\_MAX\_REJECTS} & 4 & Re-sync after this many consecutive rejected fixes\\
\code{TRAIL\_MAX} & 250 & Max trail points kept in memory\\
\bottomrule
\end{tabularx}

\section{Glossary}
\begin{description}[style=nextline,leftmargin=0pt,labelwidth=0pt,labelsep=0pt,itemindent=0pt]
  \item[BLE] Bluetooth Low Energy, the radio link from the car controller to the phone.
  \item[MQTT] A lightweight publish/subscribe messaging protocol; the broker relays messages from the phone to the engineer dashboard.
  \item[Broker] The MQTT server (e.g. HiveMQ Cloud) both sides connect to.
  \item[TLS] Encryption for the broker connection.
  \item[WebSocket] The transport the browser uses to speak MQTT (\code{ws://} / \code{wss://}).
  \item[Fuel cell (FC)] The hydrogen power source on the car.
  \item[Supercap] Supercapacitor buffer storing energy for bursts.
  \item[Trail] The line drawn behind the car on the map, here colored by motor power.
\end{description}


% #####################################################################
%  PART II - DEUTSCH
% #####################################################################
\selectlanguage{ngerman}
\guidepart{Deutsches Handbuch}

\section{Was dieses Projekt ist}
ThaiGer Track Control ist ein zweiteiliges Telemetriesystem f\"ur ein
Wasserstoff-Brennstoffzellen-Rennauto:
\begin{itemize}
  \item Eine \textbf{Android-App}, die auf einem Smartphone \emph{im Fahrzeug}
    l\"auft. Sie empf\"angt Live-Telemetrie vom Steuerger\"at des Autos \"uber
    \textbf{Bluetooth Low Energy (BLE)}, zeigt sie dem Fahrer auf einem
    Renn-Dashboard, zeichnet den Lauf auf und \textbf{leitet} die Daten \"uber das
    Internet per \textbf{MQTT} weiter.
  \item Ein \textbf{Engineer-Dashboard} --- eine einzelne HTML-Datei
    (\code{engineer\_dashboard.html}) --- das im Browser auf dem \emph{Laptop des
    Ingenieurs} in der Box l\"auft. Es abonniert denselben MQTT-Stream und zeigt
    Betriebsdaten, eine Live-Karte mit leistungsgef\"arbter GPS-Spur und
    Analyse-Diagramme.
\end{itemize}
Beide H\"alften kommunizieren nie direkt; sie treffen sich am \textbf{MQTT-Broker}
in der Cloud.

\section{System\"uberblick}

\begin{lstlisting}
   +------------+   BLE    +------------------+   MQTT/TLS   +----------------+
   | Fahrzeug-  | -------> | Android-App      | -----------> | MQTT-Broker    |
   | Steuerg.   | *A28.3*  | (Handy im Auto)  | JSON-Frames  | (HiveMQ Cloud) |
   +------------+          +------------------+              +-------+--------+
                                  |                                  | WebSocket/TLS
                                  | GPS (Handy LocationManager)      v
                                  v                          +----------------+
                          thaiger/{car}/gps                  | Engineer-      |
                          thaiger/{car}/telemetry            | Dashboard HTML |
                                                             +----------------+
\end{lstlisting}
\begin{itemize}
  \item Zwei MQTT-Topics pro Fahrzeug: \code{thaiger/\{car\}/telemetry} und
    \code{thaiger/\{car\}/gps}.
  \item \code{\{car\}} ist die Fahrzeug-ID: \code{thaiger7} oder \code{bengalo}.
  \item Das Handy ver\"offentlicht (publish); das Engineer-Dashboard abonniert
    (subscribe). Beliebig viele Laptops k\"onnen dasselbe Fahrzeug gleichzeitig
    beobachten.
\end{itemize}

\section{Die Android-App}

\subsection{Installation \& erster Start}
Die App ist ein Standard-Android-Projekt (Gradle). Debug-APK bauen:

\begin{lstlisting}
./gradlew assembleDebug
\end{lstlisting}

\begin{tabularx}{\linewidth}{@{}lL@{}}
\toprule
\textbf{Punkt} & \textbf{Wert}\\
\midrule
Minimum-Android & 9.0 (API 28)\\
Ziel-Android & 14 (API 34)\\
Sprache & Java 11\\
App-ID & \code{com.thaiger.h2racing}\\
\bottomrule
\end{tabularx}

\medskip
Installiere die APK auf dem Handy, das im Auto bleibt. Beim ersten Start \"offnet
die App den Bildschirm \textbf{Fahrzeugauswahl}.
\begin{tgnote}
\textbf{Tipp.} Montiere das Handy im Querformat so, dass der Fahrer es sieht. Das
Dashboard h\"alt den Bildschirm w\"ahrend eines Laufs wach (Einstellungen
\arrow{} \emph{Bildschirm anlassen}).
\end{tgnote}

\subsection{Berechtigungen}
Die App fragt Berechtigungen auf dem \textbf{Verbinden}-Bildschirm ab, bevor
Hardware genutzt wird. Erteile alle f\"ur volle Funktionalit\"at.

\begin{tabularx}{\linewidth}{@{}lL@{}}
\toprule
\textbf{Berechtigung} & \textbf{Wozu sie n\"otig ist}\\
\midrule
Bluetooth Connect / Scan (Android 12+) & BLE-Modul des Autos koppeln \& ansprechen\\
Bluetooth / Bluetooth Admin (Android $\le$ 11) & Dasselbe auf \"alterem Android\\
Standort (genau) & Von Android f\"ur BLE-Scan \emph{und} f\"ur GPS verlangt\\
Standort (grob) & GPS-Fallback\\
Internet & MQTT-Relay zum Ingenieur-Laptop\\
Wake Lock & Bildschirm w\"ahrend eines Laufs anlassen\\
Vibration & Haptisches Feedback bei Warnungen\\
\bottomrule
\end{tabularx}

\medskip
Verweigerst du Bluetooth, kannst du trotzdem den \textbf{Demo-Modus}
(synthetische Daten) nutzen.

\subsection{Bildschirm 1 --- Fahrzeugauswahl}
\begin{itemize}
  \item Zwei Karten: \textbf{Thaiger 7} und \textbf{Bengalo}. Tippe eine Karte an,
    um sie auszuw\"ahlen (blauer Rahmen markiert die Auswahl). Thaiger 7 ist
    Standard.
  \item Das gew\"ahlte Profil steuert Warnschwellen und Maximalleistungs-Skala.
  \item Tippe oben auf \textbf{Settings}, um vor einem Lauf die Einstellungen zu
    \"offnen.
  \item Tippe auf \textbf{Connect to \ldots}, um zum Verbinden-Bildschirm zu gehen.
\end{itemize}

\subsection{Bildschirm 2 --- Verbinden (Bluetooth / Demo-Modus)}
Dieser Bildschirm koppelt das Handy mit dem Auto und startet den Datenstrom.
\begin{enumerate}
  \item Die App pr\"uft Berechtigungen und ob Bluetooth an ist.
  \item Ein Dialog listet deine \textbf{gekoppelten Bluetooth-Ger\"ate} plus einen
    letzten Eintrag \textbf{Demo-Modus}.
  \item W\"ahle das Modul deines Autos (z.\,B. ein Ger\"at mit Namen HM-10 / HC-05
    / ESP32) \textbf{oder} den \textbf{Demo-Modus} f\"ur einen synthetischen
    Datenstrom ohne Hardware.
  \item \textbf{Pair Neu} \"offnet die Android-Bluetooth-Einstellungen, falls dein
    Modul noch nicht gekoppelt ist. \textbf{Cancel} kehrt zur Fahrzeugauswahl
    zur\"uck.
\end{enumerate}
Meldet das Modul \textbf{Connected}, wechselt die App automatisch zum
Live-Dashboard. Ist das MQTT-Relay aktiviert, startet es hier ebenfalls --- der
Ingenieur-Laptop empf\"angt Daten, sobald der Lauf beginnt.
\begin{tgnote}
Verbindungsabbr\"uche werden automatisch mit exponentiellem Backoff behandelt
(1 \arrow{} 2 \arrow{} 4 \arrow{} 8 \arrow{} 16 s). Du siehst
\emph{Reconnecting\ldots} --- kein Eingreifen n\"otig.
\end{tgnote}

\subsection{Bildschirm 3 --- Live-Dashboard}
Das Dashboard ist die Hauptansicht des Fahrers (Querformat, Bildschirm bleibt
wach). Es zeigt laufend den neuesten Wert jedes Telemetriefelds.
\begin{itemize}
  \item \textbf{Geschwindigkeit} --- die gro\ss{}e Hero-Zahl (km/h). Ist
    \emph{Geschwindigkeits-Farbwarnung} an, wird die Zahl \textbf{rot}, wenn du
    unter die Mindest-Zielgeschwindigkeit f\"allst. Ein Indikator zeigt \ok{} on
    target / \nope{} too slow.
  \item \textbf{Runden-Overlay} (eine Zeile \"uber der unteren Leiste):
    \code{PREV 1:18} \midsep{} \code{LAP 3} \midsep{} \code{CURR 0:09} ---
    vorherige Rundenzeit, aktuelle Rundennummer und laufende Rundenzeit.
  \item \textbf{Elektrik-Block} --- Spannung \& Strom von Brennstoffzelle /
    Supercap / Motor, Eigenverbrauch, Zellspannungsdifferenz.
  \item \textbf{Thermik \& Effizienz} --- BZ-Temperatur, Luftpumpen-Duty, BZ- und
    System-Wirkungsgrad, Energien, optimale Geschwindigkeit, Fahrhinweis.
  \item \textbf{Warn-Banner} --- erscheint, wenn eine Schwelle \"uberschritten wird
    (z.\,B. BZ-Temperatur oder Zellspannungsdifferenz zu hoch). Mit
    \emph{Vibration} / \emph{Warnton} an gibt es zus\"atzlich Feedback. Tippe auf
    \textbf{Dismiss}, um das Banner auszublenden.
\end{itemize}
\begin{tgnote}
\textbf{Einen Lauf beenden:} \textbf{lange auf die Geschwindigkeitszahl
dr\"ucken.} Das stoppt Bluetooth und das MQTT-Relay, schlie\ss{}t die Statistik ab
und \"offnet die \textbf{Run-Zusammenfassung}.
\end{tgnote}

\subsection{Bildschirm 4 --- Run-Zusammenfassung \& CSV-Export}
Nach einem Lauf siehst du aggregierte Ergebnisse: Laufdauer, Distanz,
Durchschnittsgeschwindigkeit, verbrauchte Energie, Spitzen-Motorleistung, maximale
BZ-Temperatur und Anzahl der Warnungen, plus ein \textbf{Leistungsdiagramm}.
\begin{itemize}
  \item \textbf{Export CSV} --- schreibt jedes aufgezeichnete Frame in eine
    CSV-Datei und \"offnet das Teilen-Men\"u. Die Datei hei\ss{}t
    \code{thaiger\_\{car\}\_\allowbreak\{YYYYMMDD\_HHMM\}.csv} und hat 21 Spalten. Ohne Daten
    erscheint \emph{\glqq No data to export\grqq{}}.
  \item \textbf{New Run} --- direkt zur\"uck zum Verbinden.
  \item \textbf{Back to Car Select} --- zur\"uck zum Anfang.
\end{itemize}

\subsection{Bildschirm 5 --- Einstellungen}
Erreichbar \"uber \textbf{Settings} in der Fahrzeugauswahl. Alle Werte werden
sofort gespeichert.

\subsubsection{Anzeige}

\begin{tabularx}{\linewidth}{@{}llL@{}}
\toprule
\textbf{Einstellung} & \textbf{Standard} & \textbf{Bedeutung}\\
\midrule
Bildschirm anlassen & An & Wake-Lock halten, damit der Bildschirm nie schl\"aft\\
Geschwindigkeits-Farbwarnung & An & Geschwindigkeit rot f\"arben, wenn unter der Mindestgeschwindigkeit\\
Aktualisierungsrate & 50 ms & Wie oft die UI neu zeichnet (20--500 ms)\\
\bottomrule
\end{tabularx}

\subsubsection{Warnungen}

\begin{tabularx}{\linewidth}{@{}llL@{}}
\toprule
\textbf{Einstellung} & \textbf{Standard} & \textbf{Bedeutung}\\
\midrule
Vibration & An & Haptischer Impuls bei einer Warnung\\
Warnton & Aus & Ton bei einer Warnung abspielen\\
\bottomrule
\end{tabularx}

\subsubsection{Schwellwerte (pro Fahrzeug)}
Separat f\"ur Thaiger 7 und Bengalo gespeichert:

\begin{tabularx}{\linewidth}{@{}llL@{}}
\toprule
\textbf{Einstellung} & \textbf{Bereich} & \textbf{Standard (Thaiger 7 / Bengalo)}\\
\midrule
BZ-Temperatur max & 30--120 \degC{} & 70 \degC{} / 65 \degC{}\\
Zellspannungsdifferenz max & 10--500 mV & 50 mV / 50 mV\\
\bottomrule
\end{tabularx}

\subsection{Bildschirm 6 --- MQTT-Einstellungen}
Unterbildschirm der Einstellungen, der das Relay zum Ingenieur-Laptop
konfiguriert. Das Relay l\"auft nur, wenn es \emph{vor} dem Start aktiviert ist.

\begin{tabularx}{\linewidth}{@{}llL@{}}
\toprule
\textbf{Feld} & \textbf{Standard} & \textbf{Hinweise}\\
\midrule
Relay aktivieren & Aus & Hauptschalter --- einschalten, um Telemetrie zu senden\\
TLS verwenden & An & Verschl\"usselte Verbindung (von HiveMQ Cloud verlangt)\\
Broker-Host & --- & z.\,B. \code{abc123.s1.eu.hivemq.cloud}\\
Port & 8883 & TLS-MQTT-Port (8883 typisch; 1883 unverschl\"usselt)\\
Benutzername & --- & Broker-Benutzername\\
Passwort & --- & Broker-Passwort (maskiert)\\
Publish-Rate & 5 Hz & Telemetrie-Senderate, 1--10 Hz\\
\bottomrule
\end{tabularx}
\begin{tgnote}
Das Handy sendet \"uber \textbf{natives MQTT (TLS)}. Das Engineer-Dashboard
verbindet sich mit demselben Broker \"uber \textbf{MQTT-\"uber-WebSocket} --- der
Broker muss also auch einen WebSocket-Listener anbieten (HiveMQ Cloud: Port 8884).
\end{tgnote}

\section{Das Engineer-Dashboard}
\code{engineer\_dashboard.html} ist eine einzelne, eigenst\"andige Datei. Sie
braucht keinen Build-Schritt und keinen Server --- nur einen modernen Browser und
Internetzugang (Kartenkacheln, Schriften, MQTT-/Diagramm-Bibliotheken vom CDN).

\subsection{\"Offnen \& verbinden}
\begin{enumerate}
  \item \"Offne \code{engineer\_dashboard.html} in Chrome/Edge/Firefox.
  \item Der Dialog \textbf{Connection Settings} \"offnet sich beim ersten Mal
    automatisch (oder klicke oben rechts auf das Zahnrad).
  \item Trage ein: \textbf{Car} (\code{Thaiger 7} oder \code{Bengalo}, muss zum
    Auto passen); \textbf{Broker URL (WebSocket)}, z.\,B.
    \code{wss://abc123.s1.eu.hivemq.cloud:\allowbreak8884/mqtt} (\code{wss://} f\"ur TLS oder
    \code{ws://host:port} lokal); \textbf{Username} / \textbf{Password}.
  \item Klicke \textbf{Connect}. Die Einstellungen werden im Browser gespeichert
    (\code{localStorage}); n\"achstes Mal automatische Verbindung.
\end{enumerate}
Der Status-Punkt oben rechts: grau getrennt, gelb verbindet, gr\"un verbunden. Drei
Tabs wechseln zwischen \textbf{Operations}, \textbf{Live Tracking} und
\textbf{Analytics}.

\subsection{Ansicht 1 --- Operations}
Gesamt\"uberblick: eine \textbf{Hero-Geschwindigkeitskarte} (blau bei/\"uber der
optimalen Geschwindigkeit, rot darunter) mit Durchschnittsgeschwindigkeit,
Rundenzahl, Laufzeit, Zielrunde und Distanz; ein \textbf{Elektrik}-Raster; ein
\textbf{Thermik \& Effizienz}-Raster; und eine pulsierende \textbf{LIVE}-Pille.

\subsection{Ansicht 2 --- Live-Tracking \& Karte}
Eine bildschirmf\"ullende dunkle \textbf{OpenStreetMap} (Leaflet), die dem Auto
folgt.
\begin{itemize}
  \item \textbf{Leistungsgef\"arbte Spur} --- segmentweise gezeichnet und nach der
    Motorleistung gef\"arbt: \textcolor{tgreen}{gr\"un (0 W)} \arrow{}
    \textcolor{tamber}{gelb (100 W)} \arrow{} \textcolor{tred}{rot (200 W+)}. Der
    neueste Teil ist voll deckend und verblasst mit dem Alter. (200 W ist die max.
    ThaiGer-Tracking-Leistung; Konstante \code{POWER\_MAX} \"anderbar.)
  \item \textbf{Leistungslegende} (unten links) --- Farbskala plus Live-Watt-Anzeige.
  \item \textbf{Rotierender Fahrzeugmarker} zeigt in die Fahrtrichtung.
  \item \textbf{GPS-Statuspanel} (rechts) --- Fix-Qualit\"at, Genauigkeit, Breite,
    L\"ange, H\"ohe, GPS-Geschwindigkeit, Kurs, Live-Motorleistung, Rundenzahl und
    ein \textbf{Filtered}-Z\"ahler.
\end{itemize}

\subsubsection{GPS-Jitter-Bereinigung}
Jeder Fix wird in vier Schritten bereinigt: (1) Fixes mit Genauigkeit schlechter
als 75 m verwerfen; (2) \glqq Teleports\grqq{} ablehnen --- Fixes mit unm\"oglicher
Geschwindigkeit (\"uber ca.\ 162 km/h), mit Re-Sync-Sicherung; (3) Wackeln unter
1,5 m im Stand ignorieren; (4) leichte Gl\"attung. Ergebnis: saubere Spur ohne
Spr\"unge.

\subsubsection{Snapshot \& Reset pro Runde}
Z\"ahlt der Rundenz\"ahler hoch, rendert das Dashboard die Spur der beendeten Runde
automatisch zu einem PNG und l\"adt es herunter
(\code{thaiger-lapN-<Zeitstempel>.png}), leert dann die Spur. Ein gr\"uner Toast
best\"atigt jede Speicherung.

\subsection{Ansicht 3 --- Analytics}
Eine \textbf{Live-Telemetrietabelle} mit allen 22 Feldern, plus ein
\textbf{Speed}-Diagramm (km/h \"uber Zeit) und ein \textbf{Efficiency}-Diagramm
(BZ- und System-Wirkungsgrad), jeweils die letzten 60 Messwerte.

\section{Telemetrie-Protokoll \& Feldreferenz}
Das Steuerger\"at sendet ASCII-Frames \"uber BLE. Jedes Feld ist in Asterisken
eingerahmt, mit einem Einzelbuchstaben als Schl\"ussel:

\begin{lstlisting}
*A28.3**B25.0**C3**D15:35**N53.4*
\end{lstlisting}
Der Parser ist stream-tolerant und beh\"alt stets den letzten bekannten Wert jedes
Felds.

\begin{tabularx}{\linewidth}{@{}clL@{}}
\toprule
\textbf{Schl.} & \textbf{Feld} & \textbf{Einheit}\\
\midrule
A & Geschwindigkeit & km/h\\
B & Durchschnittsgeschwindigkeit & km/h\\
C & Rundenzahl & ---\\
D & Gesamtlaufzeit & mm:ss\\
E & Zielrundenzeit & mm:ss\\
F & Optimale Geschwindigkeit & km/h\\
G & Brennstoffzellen-Spannung & V\\
H & Supercap-Spannung & V\\
I & Motorspannung & V\\
J & Brennstoffzellen-Strom & A\\
K & Supercap-Strom & A\\
L & Motorstrom & A\\
M & Eigenverbrauch & A\\
N & Brennstoffzellen-Temperatur & \degC{}\\
O & Luftpumpen-Duty-Cycle & \%\\
P & Fahrhinweis & ---\\
Q & Zellspannungsdifferenz & mV\\
R & Brennstoffzellen-Energie (kumuliert) & Ws\\
S & Motorenergie (kumuliert) & Ws\\
T & Brennstoffzellen-Wirkungsgrad & \%\\
U & System-Wirkungsgrad & \%\\
\bottomrule
\end{tabularx}

\medskip
Distanz, Motorleistung ($V \times A$) und Energie in Wh werden auf dem Handy
berechnet.

\section{MQTT-Topics \& Payloads}
Alle Nachrichten sind \textbf{QoS 0}. Leere/NaN-Felder werden weggelassen (Payload
ca.\ 100--200 Bytes).

\noindent\textbf{\code{thaiger/\{car\}/telemetry}} --- mit konfigurierter Rate
(Standard 5 Hz):

\begin{lstlisting}
{"ts":1736517201234,"A":28.3,"B":25.0,"C":3,"D":"15:35","G":28.3,"N":53.4,"dist":0.42}
\end{lstlisting}
\noindent\textbf{\code{thaiger/\{car\}/gps}} --- wenn ein GPS-Fix verf\"ugbar ist:

\begin{lstlisting}
{"ts":1736517201300,"lat":13.756331,"lon":100.501762,"alt":12.0,"spd":53.4,"acc":3.5,"hdg":270.0}
\end{lstlisting}

\begin{tabularx}{\linewidth}{@{}lLl@{}}
\toprule
\textbf{Feld} & \textbf{Bedeutung} & \textbf{Einheit}\\
\midrule
ts & Unix-Zeitstempel & ms\\
lat / lon & Breite / L\"ange & Grad\\
alt & H\"ohe (falls verf\"ugbar) & m\\
spd & GPS-Geschwindigkeit (falls verf\"ugbar) & km/h\\
acc & Horizontale Genauigkeit (falls verf\"ugbar) & m\\
hdg & Kurs / Peilung (falls verf\"ugbar) & Grad\\
\bottomrule
\end{tabularx}

\section{Fahrzeugprofile}

\begin{tabularx}{\linewidth}{@{}llllL@{}}
\toprule
\textbf{Profil} & \textbf{id} & \textbf{Max. Leistung} & \textbf{BZ-Temp-Limit} & \textbf{Mindestgeschw.}\\
\midrule
Thaiger 7 & \code{thaiger7} & 600 W & 70 \degC{} & 25 km/h\\
Bengalo & \code{bengalo} & 900 W & 65 \degC{} & 25 km/h\\
\bottomrule
\end{tabularx}

\medskip
Schwellwerte k\"onnen pro Auto in den \textbf{Einstellungen} \"uberschrieben
werden. Die Kartenspur des Engineer-Dashboards nutzt eine separate \textbf{200 W}-
Leistungsskala (\code{POWER\_MAX}).

\section{Fehlerbehebung}

\begin{tabularx}{\linewidth}{@{}LL@{}}
\toprule
\textbf{Symptom} & \textbf{Ursache / Abhilfe}\\
\midrule
App verbindet nicht \"uber Bluetooth & Modul zuerst in den Android-Bluetooth-Einstellungen koppeln; Stromversorgung pr\"ufen; Bluetooth- + Standort-Berechtigungen erteilen. \textbf{Pair Neu} nutzen.\\
Keine Daten, aber \glqq Connected\grqq{} & Falsches Ger\"at oder Steuerger\"at sendet nicht. Demo-Modus testen, dann Modul kontrollieren.\\
Dashboard wird nie gr\"un & Falsche WebSocket-URL (muss \code{wss://\ldots:8884/mqtt} sein, nicht \code{8883}), falsche Zugangsdaten oder kein WebSocket-Listener.\\
Dashboard gr\"un, aber keine Telemetrie & Das \textbf{Car} passt nicht zur Fahrzeug-ID, oder das Relay ist am Handy nicht aktiviert. Beide brauchen dasselbe \code{thaiger7}/\code{bengalo}.\\
Karte ist leer & Kein Internet f\"ur OSM-Kacheln oder noch keine GPS-Fixes. Die Karte zentriert erst mit einem Fix.\\
GPS-Spur springt & Die Bereinigung sollte das verhindern; bei nur groben Fixes \code{GPS\_MAX\_ACC} erh\"ohen. Der \textbf{Filtered}-Z\"ahler zeigt verworfene Fixes.\\
GPS aktualisiert nur alle paar Sekunden & Freie Sicht zum Himmel und \textbf{Standort}-Berechtigung sicherstellen; drinnen grobe Netzwerk-Fixes.\\
Kein CSV erzeugt & Lauf mit aufgezeichneten Frames abschlie\ss{}en; sonst meldet \emph{Export CSV} \glqq No data to export\grqq{}.\\
\bottomrule
\end{tabularx}

\section{Feinabstimmung}
Am Handy (Einstellungen / MQTT-Einstellungen): Bildschirm-Wach, Geschwindigkeits-
Farbe, UI-Rate, Vibration/Ton, Schwellwerte pro Auto, Broker-/Relay-Parameter.

\medskip
\noindent In \code{engineer\_dashboard.html} (Konstanten oben im Script-Block):

\begin{tabularx}{\linewidth}{@{}llL@{}}
\toprule
\textbf{Konstante} & \textbf{Standard} & \textbf{Zweck}\\
\midrule
\code{POWER\_MAX} & 200 & Watt-Wert f\"ur volles Rot der Spur/Legende\\
\code{GPS\_MAX\_ACC} & 75 & GPS-Fixes schlechter als so viele Meter verwerfen\\
\code{GPS\_MAX\_SPEED} & 45 & m/s; schnellere implizierte Geschw. = Teleport\\
\code{GPS\_MIN\_MOVE} & 1.5 & m; kleinere Bewegungen = Stillstands-Jitter\\
\code{GPS\_SMOOTH} & 0.35 & EMA-Gl\"attungsfaktor (0 = keine, 1 = roh)\\
\code{GPS\_MAX\_REJECTS} & 4 & Re-Sync nach so vielen verworfenen Fixes\\
\code{TRAIL\_MAX} & 250 & Max. Anzahl Spurpunkte im Speicher\\
\bottomrule
\end{tabularx}

\section{Glossar}
\begin{description}[style=nextline,leftmargin=0pt,labelwidth=0pt,labelsep=0pt,itemindent=0pt]
  \item[BLE] Bluetooth Low Energy, die Funkverbindung vom Steuerger\"at zum Handy.
  \item[MQTT] Leichtgewichtiges Publish/Subscribe-Protokoll; der Broker leitet Nachrichten weiter.
  \item[Broker] Der MQTT-Server (z.\,B. HiveMQ Cloud), mit dem sich beide Seiten verbinden.
  \item[TLS] Verschl\"usselung der Broker-Verbindung.
  \item[WebSocket] Der Transport, \"uber den der Browser MQTT spricht (\code{ws://} / \code{wss://}).
  \item[Brennstoffzelle (BZ)] Die Wasserstoff-Energiequelle im Auto.
  \item[Supercap] Superkondensator-Puffer f\"ur Lastspitzen.
  \item[Spur (Trail)] Die hinter dem Auto gezeichnete Linie, nach Motorleistung gef\"arbt.
\end{description}


% #####################################################################
%  PART III - CHINESE (Simplified)
% #####################################################################
\selectlanguage{english}
\guidepart{中文使用手册}

\section{本项目是什么}
ThaiGer Track Control 是一套面向氢燃料电池赛车的两端式遥测系统：
\begin{itemize}
  \item 一个运行在\emph{车内手机}上的 \textbf{安卓 App}。它通过\textbf{低功耗蓝牙
    （BLE）}接收赛车控制器的实时遥测数据，在赛车仪表盘上展示给车手，记录本次行驶，
    并通过 \textbf{MQTT} 经互联网\textbf{转发}数据。
  \item 一个 \textbf{工程师仪表盘} —— 一个单独的 HTML 文件
    （\code{engineer\_dashboard.html}）—— 运行在维修区\emph{工程师笔记本电脑}的浏览
    器中。它订阅同一条 MQTT 数据流，显示运营数据、带功率着色 GPS 轨迹的实时地图，以及
    分析图表。
\end{itemize}
两端从不直接通信；它们在云端的 \textbf{MQTT 代理服务器（broker）} 处汇合。

\section{系统总览}

\begin{lstlisting}
   +------------+   BLE    +------------------+   MQTT/TLS   +----------------+
   | Car        | -------> | Android app      | -----------> | MQTT broker    |
   | controller | *A28.3*  | (phone in car)   | JSON frames  | (HiveMQ Cloud) |
   +------------+          +------------------+              +-------+--------+
                                  |                                  | WebSocket/TLS
                                  | GPS (phone LocationManager)      v
                                  v                          +----------------+
                          thaiger/{car}/gps                  | Engineer       |
                          thaiger/{car}/telemetry            | Dashboard HTML |
                                                             +----------------+
\end{lstlisting}
\begin{itemize}
  \item 每辆车有两个 MQTT 主题：\code{thaiger/\{car\}/telemetry} 与
    \code{thaiger/\{car\}/gps}。
  \item \code{\{car\}} 是车辆 ID：\code{thaiger7} 或 \code{bengalo}。
  \item 手机负责发布（publish），工程师仪表盘负责订阅（subscribe）。任意数量的笔记本
    电脑可同时观看同一辆车。
\end{itemize}

\section{安卓 App}

\subsection{安装与首次启动}
本 App 是标准的安卓项目（Gradle）。构建调试版 APK：

\begin{lstlisting}
./gradlew assembleDebug
\end{lstlisting}

\begin{tabularx}{\linewidth}{@{}lL@{}}
\toprule
\textbf{项目} & \textbf{取值}\\
\midrule
最低安卓版本 & 9.0（API 28）\\
目标安卓版本 & 14（API 34）\\
语言 & Java 11\\
应用 ID & \code{com.thaiger.h2racing}\\
\bottomrule
\end{tabularx}

\medskip
将 APK 安装到将放置在车内的手机上。首次启动时，App 打开 \textbf{车辆选择} 界面。
\begin{tgnote}
\textbf{提示。} 以横屏方式将手机固定在车手能看到的位置。行驶过程中仪表盘会保持屏幕常亮
（设置 \arrow{} \emph{保持屏幕常亮}）。
\end{tgnote}

\subsection{权限}
App 在使用任何硬件之前，会在 \textbf{连接} 界面请求权限。请全部授予以获得完整功能。

\begin{tabularx}{\linewidth}{@{}lL@{}}
\toprule
\textbf{权限} & \textbf{用途}\\
\midrule
蓝牙连接 / 扫描（安卓 12+） & 配对并与车辆 BLE 模块通信\\
蓝牙 / 蓝牙管理（安卓 $\le$ 11） & 同上，用于较旧安卓\\
定位（精确） & 安卓要求其用于 BLE 扫描\emph{以及} GPS\\
定位（粗略） & GPS 回退\\
互联网 & 向工程师笔记本电脑做 MQTT 转发\\
Wake Lock & 行驶中保持屏幕常亮\\
振动 & 报警时的触感反馈\\
\bottomrule
\end{tabularx}

\medskip
若拒绝蓝牙权限，你仍可使用 \textbf{演示模式}（合成数据）来体验 App。

\subsection{界面 1 —— 车辆选择}
\begin{itemize}
  \item 两张卡片：\textbf{Thaiger 7} 与 \textbf{Bengalo}。点按卡片即可选择（蓝色描边
    标记当前选择）。默认选中 Thaiger 7。
  \item 所选配置档决定报警阈值与最大功率刻度。
  \item 点按顶部 \textbf{Settings}，在行驶前打开 App 设置。
  \item 点按 \textbf{Connect to \ldots} 进入连接界面。
\end{itemize}

\subsection{界面 2 —— 连接（蓝牙 / 演示模式）}
此界面把手机与赛车配对并启动数据流。
\begin{enumerate}
  \item App 检查权限以及蓝牙是否开启。
  \item 弹出对话框，列出你\textbf{已配对的蓝牙设备}，并在末尾附上一项 \textbf{演示模式}。
  \item 选择你赛车的模块（例如名称含 HM-10 / HC-05 / ESP32 的设备），\textbf{或}选择
    \textbf{演示模式}以在无硬件的情况下运行合成数据流。
  \item \textbf{Pair Neu / 配对新设备} 会打开安卓蓝牙设置（若模块尚未配对）。
    \textbf{Cancel} 返回车辆选择。
\end{enumerate}
当模块报告 \textbf{Connected（已连接）} 后，App 会自动跳转到实时仪表盘。若已启用 MQTT
转发，它也会在此启动 —— 一旦行驶开始，工程师笔记本电脑便开始接收数据。
\begin{tgnote}
连接中断会以指数退避自动处理（1 \arrow{} 2 \arrow{} 4 \arrow{} 8 \arrow{} 16 秒）。
你会看到 \emph{Reconnecting\ldots}，无需任何操作。
\end{tgnote}

\subsection{界面 3 —— 实时仪表盘}
仪表盘是车手的主视图（横屏，屏幕保持常亮）。它持续刷新，显示每个遥测字段的最新值。
\begin{itemize}
  \item \textbf{速度} —— 大号主数字（km/h）。若开启了\emph{速度颜色警告}，当你低于该车
    的最低目标速度时数字变为\textbf{红色}。一个小指示器显示 \ok{} on target /
    \nope{} too slow。
  \item \textbf{圈速浮层}（底部栏上方的一行）：\code{PREV 1:18} \midsep{} \code{LAP 3}
    \midsep{} \code{CURR 0:09} —— 上一圈用时、当前圈号，以及当前圈的实时计时。
  \item \textbf{电气区块} —— 燃料电池 / 超级电容 / 电机的电压与电流、自身消耗、电芯电压差。
  \item \textbf{热与效率} —— 燃料电池温度、空气泵占空比、燃料电池与系统效率、能量、
    最优速度、驾驶提示。
  \item \textbf{报警横幅} —— 当某阈值被越过时出现（例如燃料电池温度或电芯电压差过高）。
    开启\emph{振动} / \emph{报警声} 后还会有触感 / 声音反馈。点按 \textbf{Dismiss}
    隐藏横幅。
\end{itemize}
\begin{tgnote}
\textbf{结束一次行驶：} \textbf{长按速度数字。} 这会停止蓝牙与 MQTT 转发，结算本次行驶
统计，并打开 \textbf{赛后总结}。
\end{tgnote}

\subsection{界面 4 —— 赛后总结与 CSV 导出}
行驶结束后你会看到汇总结果：行驶时长、距离、平均速度、消耗能量、电机峰值功率、燃料电池
最高温度，以及报警次数，还有整段行驶的\textbf{功率曲线图}。
\begin{itemize}
  \item \textbf{Export CSV（导出 CSV）} —— 将每一帧已记录的遥测数据写入 CSV 文件，并
    打开安卓分享面板。文件名为 \code{thaiger\_\{car\}\_\allowbreak\{YYYYMMDD\_HHMM\}.csv}，包含 21
    列。若未记录到数据，会提示 \emph{``No data to export''}。
  \item \textbf{New Run（新行驶）} —— 直接返回连接界面开始下一次行驶。
  \item \textbf{Back to Car Select（返回车辆选择）} —— 回到起点。
\end{itemize}

\subsection{界面 5 —— 设置}
通过车辆选择界面的 \textbf{Settings} 链接进入。所有取值立即保存。

\subsubsection{显示}

\begin{tabularx}{\linewidth}{@{}llL@{}}
\toprule
\textbf{设置} & \textbf{默认} & \textbf{含义}\\
\midrule
保持屏幕常亮 & 开 & 持有 wake lock，使屏幕不休眠\\
速度颜色警告 & 开 & 低于该车最低目标速度时把速度数字变红\\
刷新频率 & 50 ms & 仪表盘 UI 的重绘间隔（20--500 ms）\\
\bottomrule
\end{tabularx}

\subsubsection{报警}

\begin{tabularx}{\linewidth}{@{}llL@{}}
\toprule
\textbf{设置} & \textbf{默认} & \textbf{含义}\\
\midrule
振动 & 开 & 报警触发时的触感脉冲\\
报警声 & 关 & 报警触发时播放提示音\\
\bottomrule
\end{tabularx}

\subsubsection{阈值（每辆车独立）}
为 Thaiger 7 与 Bengalo 分别保存：

\begin{tabularx}{\linewidth}{@{}llL@{}}
\toprule
\textbf{设置} & \textbf{范围} & \textbf{默认（Thaiger 7 / Bengalo）}\\
\midrule
燃料电池温度上限 & 30--120 \degC{} & 70 \degC{} / 65 \degC{}\\
电芯电压差上限 & 10--500 mV & 50 mV / 50 mV\\
\bottomrule
\end{tabularx}

\subsection{界面 6 —— MQTT 设置}
设置的一个子界面，用于配置到工程师笔记本电脑的转发。仅当在开始行驶\emph{之前}启用，转发
才会运行。

\begin{tabularx}{\linewidth}{@{}llL@{}}
\toprule
\textbf{字段} & \textbf{默认} & \textbf{说明}\\
\midrule
启用转发 & 关 & 总开关 —— 打开后才会发布遥测\\
使用 TLS & 开 & 加密连接（HiveMQ Cloud 要求）\\
Broker 主机 & --- & 例如 \code{abc123.s1.eu.hivemq.cloud}\\
端口 & 8883 & TLS MQTT 端口（通常 8883；明文用 1883）\\
用户名 & --- & Broker 用户名\\
密码 & --- & Broker 密码（以掩码显示）\\
发布频率 & 5 Hz & 遥测发布频率，1--10 Hz\\
\bottomrule
\end{tabularx}
\begin{tgnote}
手机通过\textbf{原生 MQTT（TLS）}发布。工程师仪表盘通过 \textbf{MQTT over WebSocket}
连接到同一个 broker —— 因此 broker 还必须暴露一个 WebSocket 监听端口（HiveMQ Cloud
为端口 8884）。
\end{tgnote}

\section{工程师仪表盘}
\code{engineer\_dashboard.html} 是一个单一、自包含的文件。它无需构建步骤、无需服务器
—— 只要一个现代浏览器与互联网连接（地图瓦片、字体，以及来自 CDN 的 MQTT / 图表库）。

\subsection{打开与连接}
\begin{enumerate}
  \item 在 Chrome/Edge/Firefox 中打开 \code{engineer\_dashboard.html}。
  \item 首次使用时 \textbf{Connection Settings（连接设置）} 对话框会自动弹出（或点击右
    上角的齿轮按钮）。
  \item 填写：\textbf{Car}（\code{Thaiger 7} 或 \code{Bengalo}，必须与手机发布时所用
    的车辆一致）；\textbf{Broker URL (WebSocket)}，例如
    \code{wss://abc123.s1.eu.hivemq.cloud:\allowbreak8884/mqtt}（TLS 用 \code{wss://}，本地明文用
    \code{ws://host:port}）；\textbf{Username} / \textbf{Password}。
  \item 点击 \textbf{Connect}。设置会保存在浏览器中（\code{localStorage}），下次打开会
    自动重连。
\end{enumerate}
右上角的状态圆点：灰色未连接、琥珀色连接中、绿色已连接。顶部三个标签页在
\textbf{Operations}、\textbf{Live Tracking}、\textbf{Analytics} 间切换。

\subsection{视图 1 —— Operations 运营总览}
一眼总览：一张\textbf{主速度卡片}（达到/高于最优速度时为蓝色，低于则为红色），含平均
速度、圈数、行驶时间、目标圈速与距离；一个\textbf{电气}网格；一个\textbf{热与效率}
网格；以及一个跳动的 \textbf{LIVE} 标记。

\subsection{视图 2 —— Live Tracking 实时追踪与地图}
一张全屏暗色 \textbf{OpenStreetMap}（Leaflet），跟随赛车。
\begin{itemize}
  \item \textbf{功率着色轨迹} —— 按段绘制，并按该时刻的电机功率着色：
    \textcolor{tgreen}{绿色（0 W）} \arrow{} \textcolor{tamber}{琥珀（100 W）}
    \arrow{} \textcolor{tred}{红色（200 W+）}。最新部分完全不透明，随时间淡出。
    （200 W 为 ThaiGer 追踪时的最大功率；常量 \code{POWER\_MAX} 可修改。）
  \item \textbf{功率图例}（左下角）—— 颜色刻度加上实时瓦数读数。
  \item \textbf{旋转车辆标记} 指向行驶方向。
  \item \textbf{GPS 状态面板}（右侧）—— 定位质量、精度、纬度、经度、高度、GPS 速度、
    航向、实时电机功率、圈数，以及一个 \textbf{Filtered} 计数器。
\end{itemize}

\subsubsection{GPS 抖动清理}
每个定位在进入地图前都会经过四步清理：（1）剔除精度差于 75 m 的定位；（2）拒绝“瞬移”
—— 隐含速度不可能（超过约 162 km/h）的定位，并带有重新同步保护；（3）静止时忽略小于
1.5 m 的抖动；（4）对残余噪声做轻度平滑。最终得到无随机跳变的干净轨迹。

\subsubsection{每圈快照与重置}
每当赛车圈数计数器递增时，仪表盘会自动把刚结束那一圈的轨迹渲染为 PNG 图片并下载
（\code{thaiger-lapN-<时间戳>.png}），随后清空轨迹。每次保存都会有一个绿色提示条确认。

\subsection{视图 3 —— Analytics 分析}
一张\textbf{实时遥测表格}，列出全部 22 个字段及其当前值；一张 \textbf{Speed（速度）}
图（km/h 随时间）与一张 \textbf{Efficiency（效率）} 图（燃料电池与系统效率随时间），
各显示最近约 60 个采样点。

\section{遥测协议与字段参考}
赛车控制器通过 BLE 发送 ASCII 帧。每个字段以星号包裹，并以单个字母作为键：

\begin{lstlisting}
*A28.3**B25.0**C3**D15:35**N53.4*
\end{lstlisting}
解析器对数据流容错，并始终保留每个字段的最近已知值。

\begin{tabularx}{\linewidth}{@{}clL@{}}
\toprule
\textbf{键} & \textbf{字段} & \textbf{单位}\\
\midrule
A & 速度 & km/h\\
B & 平均速度 & km/h\\
C & 圈数 & ---\\
D & 总行驶时间 & mm:ss\\
E & 目标圈速 & mm:ss\\
F & 最优速度 & km/h\\
G & 燃料电池电压 & V\\
H & 超级电容电压 & V\\
I & 电机电压 & V\\
J & 燃料电池电流 & A\\
K & 超级电容电流 & A\\
L & 电机电流 & A\\
M & 自身消耗 & A\\
N & 燃料电池温度 & \degC{}\\
O & 空气泵占空比 & \%\\
P & 驾驶提示 & ---\\
Q & 电芯电压差 & mV\\
R & 燃料电池累计能量 & Ws\\
S & 电机累计能量 & Ws\\
T & 燃料电池效率 & \%\\
U & 系统效率 & \%\\
\bottomrule
\end{tabularx}

\medskip
距离、电机功率（$V \times A$）以及以 Wh 表示的能量均在手机端推算。

\section{MQTT 主题与负载}
所有消息均为 \textbf{QoS 0}。空值/NaN 字段会被省略（负载约 100--200 字节）。

\noindent\textbf{\code{thaiger/\{car\}/telemetry}} —— 以配置频率发布（默认 5 Hz）：

\begin{lstlisting}
{"ts":1736517201234,"A":28.3,"B":25.0,"C":3,"D":"15:35","G":28.3,"N":53.4,"dist":0.42}
\end{lstlisting}
\noindent\textbf{\code{thaiger/\{car\}/gps}} —— 在有 GPS 定位时发布：

\begin{lstlisting}
{"ts":1736517201300,"lat":13.756331,"lon":100.501762,"alt":12.0,"spd":53.4,"acc":3.5,"hdg":270.0}
\end{lstlisting}

\begin{tabularx}{\linewidth}{@{}lLl@{}}
\toprule
\textbf{字段} & \textbf{含义} & \textbf{单位}\\
\midrule
ts & Unix 时间戳 & ms\\
lat / lon & 纬度 / 经度 & 度\\
alt & 高度（若可用） & m\\
spd & GPS 速度（若可用） & km/h\\
acc & 水平精度（若可用） & m\\
hdg & 航向 / 方位（若可用） & 度\\
\bottomrule
\end{tabularx}

\section{车辆配置档}

\begin{tabularx}{\linewidth}{@{}llllL@{}}
\toprule
\textbf{配置档} & \textbf{id} & \textbf{最大功率} & \textbf{燃料电池温度上限} & \textbf{最低目标速度}\\
\midrule
Thaiger 7 & \code{thaiger7} & 600 W & 70 \degC{} & 25 km/h\\
Bengalo & \code{bengalo} & 900 W & 65 \degC{} & 25 km/h\\
\bottomrule
\end{tabularx}

\medskip
阈值可在 \textbf{设置} 中按车覆盖。工程师仪表盘的地图轨迹使用独立的 \textbf{200 W}
功率刻度（\code{POWER\_MAX}）。

\section{故障排查}

\begin{tabularx}{\linewidth}{@{}LL@{}}
\toprule
\textbf{现象} & \textbf{原因 / 解决}\\
\midrule
App 无法通过蓝牙连接 & 先在安卓蓝牙设置中配对模块；检查其已通电；授予蓝牙 + 定位权限。使用 \textbf{Pair Neu / 配对新设备}。\\
显示“Connected”却无数据 & 选错了设备，或控制器未发送。先用演示模式确认 App，再检查模块。\\
仪表盘始终不变绿 & WebSocket URL 错误（必须是 \code{wss://\ldots:8884/mqtt}，而非 \code{8883}）、凭据错误，或 broker 没有 WebSocket 监听端口。\\
仪表盘变绿但无遥测 & 仪表盘中的 \textbf{Car} 与手机车辆 ID 不一致，或手机未启用转发。两端必须相同 \code{thaiger7}/\code{bengalo}。\\
地图空白 & 没有互联网加载 OSM 瓦片，或尚无 GPS 定位。地图只有在收到定位后才会居中。\\
GPS 轨迹乱跳 & 清理机制本应避免；若设备只有粗糙定位，调高 \code{GPS\_MAX\_ACC}。\textbf{Filtered} 计数器显示被剔除的定位数。\\
GPS 每隔几秒才更新一次 & 确保开阔天空视野并授予\textbf{定位}权限；室内会回退到粗糙的网络定位。\\
未生成 CSV & 必须完成一次有记录帧的行驶；否则 \emph{Export CSV} 会提示 “No data to export”。\\
\bottomrule
\end{tabularx}

\section{调参与高级设置}
在手机上（设置 / MQTT 设置）：屏幕常亮、速度颜色、UI 刷新频率、振动/声音、每车阈值，以
及所有 broker/转发参数。

\medskip
\noindent 在 \code{engineer\_dashboard.html} 中（脚本块顶部的常量）：

\begin{tabularx}{\linewidth}{@{}llL@{}}
\toprule
\textbf{常量} & \textbf{默认} & \textbf{用途}\\
\midrule
\code{POWER\_MAX} & 200 & 映射为轨迹/图例满红的瓦数值\\
\code{GPS\_MAX\_ACC} & 75 & 剔除精度差于此米数的 GPS 定位\\
\code{GPS\_MAX\_SPEED} & 45 & m/s；隐含速度更快即视为瞬移\\
\code{GPS\_MIN\_MOVE} & 1.5 & m；更小的移动视为静止抖动\\
\code{GPS\_SMOOTH} & 0.35 & EMA 平滑系数（0 = 不平滑，1 = 原始）\\
\code{GPS\_MAX\_REJECTS} & 4 & 连续剔除这么多次后重新同步\\
\code{TRAIL\_MAX} & 250 & 内存中保留的最大轨迹点数\\
\bottomrule
\end{tabularx}

\section{术语表}
\begin{description}[style=nextline,leftmargin=0pt,labelwidth=0pt,labelsep=0pt,itemindent=0pt]
  \item[BLE] 低功耗蓝牙，赛车控制器到手机的无线链路。
  \item[MQTT] 一种轻量级的发布/订阅消息协议；broker 把消息从手机转发到工程师仪表盘。
  \item[Broker（代理服务器）] 双方都连接的 MQTT 服务器（例如 HiveMQ Cloud）。
  \item[TLS] broker 连接的加密。
  \item[WebSocket] 浏览器用于承载 MQTT 的传输方式（\code{ws://} / \code{wss://}）。
  \item[燃料电池（FC）] 赛车上的氢能源。
  \item[超级电容（Supercap）] 用于应对瞬时大功率的储能缓冲。
  \item[轨迹（Trail）] 在地图上车后绘制的线，按电机功率着色。
\end{description}


% #####################################################################
%  PART IV - FRANÇAIS
% #####################################################################
\selectlanguage{french}
\guidepart{Guide en français}

\section{Présentation du projet}
ThaiGer Track Control est un système de télémétrie en deux parties pour une
voiture de course à pile à combustible à hydrogène :
\begin{itemize}
  \item Une \textbf{application Android} qui tourne sur un téléphone \emph{à bord
    de la voiture}. Elle reçoit la télémétrie en direct du calculateur de la
    voiture par \textbf{Bluetooth Low Energy (BLE)}, l'affiche sur un tableau de
    bord de course pour le pilote, enregistre la session et \textbf{retransmet} les
    données par internet via \textbf{MQTT}.
  \item Un \textbf{tableau de bord ingénieur} --- un unique fichier HTML
    (\code{engineer\_dashboard.html}) --- qui tourne dans un navigateur sur le
    \emph{portable de l'ingénieur} au stand. Il s'abonne au même flux MQTT et
    affiche les données d'exploitation, une carte en direct avec une trace GPS
    colorée par la puissance, et des graphiques d'analyse.
\end{itemize}
Les deux moitiés ne communiquent jamais directement ; elles se rejoignent au
\textbf{courtier MQTT} dans le cloud.

\section{Vue d'ensemble du système}

\begin{lstlisting}
   +------------+   BLE    +------------------+   MQTT/TLS   +----------------+
   | Car        | -------> | Android app      | -----------> | MQTT broker    |
   | controller | *A28.3*  | (phone in car)   | JSON frames  | (HiveMQ Cloud) |
   +------------+          +------------------+              +-------+--------+
                                  |                                  | WebSocket/TLS
                                  | GPS (phone LocationManager)      v
                                  v                          +----------------+
                          thaiger/{car}/gps                  | Engineer       |
                          thaiger/{car}/telemetry            | Dashboard HTML |
                                                             +----------------+
\end{lstlisting}
\begin{itemize}
  \item Deux sujets MQTT par voiture : \code{thaiger/\{car\}/telemetry} et
    \code{thaiger/\{car\}/gps}.
  \item \code{\{car\}} est l'identifiant du véhicule : \code{thaiger7} ou
    \code{bengalo}.
  \item Le téléphone publie ; le tableau de bord ingénieur s'abonne. Un nombre
    quelconque de portables peut suivre la même voiture en même temps.
\end{itemize}

\section{L'application Android}

\subsection{Installation \& premier démarrage}
L'application est un projet Android standard (Gradle). Pour construire un APK de
débogage :

\begin{lstlisting}
./gradlew assembleDebug
\end{lstlisting}

\begin{tabularx}{\linewidth}{@{}lL@{}}
\toprule
\textbf{Élément} & \textbf{Valeur}\\
\midrule
Android minimum & 9.0 (API 28)\\
Android cible & 14 (API 34)\\
Langage & Java 11\\
Identifiant & \code{com.thaiger.h2racing}\\
\bottomrule
\end{tabularx}

\medskip
Installez l'APK sur le téléphone qui reste dans la voiture. Au premier lancement,
l'application ouvre l'écran \textbf{Sélection du véhicule}.
\begin{tgnote}
\textbf{Astuce.} Montez le téléphone en paysage, visible par le pilote. Le tableau
de bord garde l'écran allumé pendant une session (Réglages \arrow{} \emph{Garder
l'écran allumé}).
\end{tgnote}

\subsection{Autorisations}
L'application demande les autorisations sur l'écran \textbf{Connexion}, avant
d'utiliser le matériel. Accordez-les toutes pour un fonctionnement complet.

\begin{tabularx}{\linewidth}{@{}lL@{}}
\toprule
\textbf{Autorisation} & \textbf{Pourquoi elle est nécessaire}\\
\midrule
Bluetooth Connect / Scan (Android 12+) & Appairer \& dialoguer avec le module BLE\\
Bluetooth / Bluetooth Admin (Android $\le$ 11) & Idem, sur Android plus ancien\\
Localisation (précise) & Exigée par Android pour le scan BLE \emph{et} le GPS\\
Localisation (approximative) & Repli GPS\\
Internet & Relais MQTT vers le portable ingénieur\\
Wake Lock & Garder l'écran allumé pendant une session\\
Vibreur & Retour haptique pour les alertes\\
\bottomrule
\end{tabularx}

\medskip
Si vous refusez le Bluetooth, vous pouvez toujours utiliser le \textbf{mode démo}
(données synthétiques).

\subsection{Écran 1 --- Sélection du véhicule}
\begin{itemize}
  \item Deux cartes : \textbf{Thaiger 7} et \textbf{Bengalo}. Touchez une carte
    pour la sélectionner (un contour bleu marque la sélection). Thaiger 7 par
    défaut.
  \item Le profil choisi pilote les seuils d'alerte et l'échelle de puissance max.
  \item Touchez \textbf{Settings} (en haut) pour ouvrir les réglages avant une
    session.
  \item Touchez \textbf{Connect to \ldots} pour passer à l'écran Connexion.
\end{itemize}

\subsection{Écran 2 --- Connexion (Bluetooth / mode démo)}
Cet écran appaire le téléphone à la voiture et démarre le flux de données.
\begin{enumerate}
  \item L'application vérifie les autorisations et que le Bluetooth est activé.
  \item Une boîte de dialogue liste vos \textbf{appareils Bluetooth appairés} plus
    une dernière entrée \textbf{mode démo}.
  \item Choisissez le module de votre voiture (p.\,ex. un appareil nommé HM-10 /
    HC-05 / ESP32), \textbf{ou} choisissez le \textbf{mode démo} pour un flux
    synthétique sans matériel.
  \item \textbf{Pair New} ouvre les réglages Bluetooth d'Android si le module n'est
    pas encore appairé. \textbf{Cancel} revient à la sélection du véhicule.
\end{enumerate}
Quand le module indique \textbf{Connected}, l'application passe automatiquement au
tableau de bord en direct. Si le relais MQTT est activé, il démarre aussi ici : le
portable ingénieur reçoit les données dès le début de la session.
\begin{tgnote}
Les coupures de connexion sont gérées automatiquement avec un délai exponentiel
(1 \arrow{} 2 \arrow{} 4 \arrow{} 8 \arrow{} 16 s). Vous verrez
\emph{Reconnecting\ldots} --- aucune action requise.
\end{tgnote}

\subsection{Écran 3 --- Tableau de bord en direct}
Le tableau de bord est la vue principale du pilote (paysage, écran allumé). Il
affiche en continu la dernière valeur de chaque champ de télémétrie.
\begin{itemize}
  \item \textbf{Vitesse} --- le grand chiffre principal (km/h). Si l'\emph{alerte
    couleur de vitesse} est active, le chiffre devient \textbf{rouge} sous la
    vitesse cible minimale. Un indicateur montre \ok{} on target / \nope{} too slow.
  \item \textbf{Bandeau de tours} (une ligne au-dessus de la barre du bas) :
    \code{PREV 1:18} \midsep{} \code{LAP 3} \midsep{} \code{CURR 0:09} --- temps du
    tour précédent, numéro du tour courant et temps courant.
  \item \textbf{Bloc électrique} --- tension \& courant pile / supercondensateur /
    moteur, consommation propre, écart de tension de cellule.
  \item \textbf{Thermique \& rendement} --- température pile, rapport cyclique de la
    pompe à air, rendement pile et système, énergies, vitesse optimale, conseil de
    conduite.
  \item \textbf{Bandeau d'alerte} --- apparaît quand un seuil est franchi (p.\,ex.
    température pile ou écart de tension trop élevé). Avec \emph{Vibreur} /
    \emph{Son d'alerte}, retour haptique / sonore en plus. Touchez \textbf{Dismiss}
    pour masquer le bandeau.
\end{itemize}
\begin{tgnote}
\textbf{Terminer une session :} \textbf{appui long sur le chiffre de vitesse.} Cela
arrête le Bluetooth et le relais MQTT, finalise les statistiques et ouvre le
\textbf{résumé de session}.
\end{tgnote}

\subsection{Écran 4 --- Résumé de session \& export CSV}
Après une session, vous voyez des résultats agrégés : durée, distance, vitesse
moyenne, énergie consommée, puissance moteur maximale, température pile maximale et
nombre d'alertes, plus un \textbf{graphique de puissance}.
\begin{itemize}
  \item \textbf{Export CSV} --- écrit chaque trame enregistrée dans un fichier CSV
    et ouvre le panneau de partage Android. Le fichier s'appelle
    \code{thaiger\_\{car\}\_\allowbreak\{YYYYMMDD\_HHMM\}.csv} et contient 21
    colonnes. Sans données, le message \emph{« No data to export »} apparaît.
  \item \textbf{New Run} --- retour direct à la Connexion.
  \item \textbf{Back to Car Select} --- retour au début.
\end{itemize}

\subsection{Écran 5 --- Réglages}
Accessible via \textbf{Settings} sur la sélection du véhicule. Tout est enregistré
immédiatement.

\subsubsection{Affichage}

\begin{tabularx}{\linewidth}{@{}llL@{}}
\toprule
\textbf{Réglage} & \textbf{Défaut} & \textbf{Signification}\\
\midrule
Garder l'écran allumé & Activé & Maintient un wake lock pour que l'écran ne s'éteigne pas\\
Alerte couleur de vitesse & Activé & Vitesse en rouge sous la vitesse cible minimale\\
Fréquence de rafraîchissement & 50 ms & Cadence de redessin de l'IU (20--500 ms)\\
\bottomrule
\end{tabularx}

\subsubsection{Alertes}

\begin{tabularx}{\linewidth}{@{}llL@{}}
\toprule
\textbf{Réglage} & \textbf{Défaut} & \textbf{Signification}\\
\midrule
Vibreur & Activé & Impulsion haptique à chaque alerte\\
Son d'alerte & Désactivé & Joue un son à chaque alerte\\
\bottomrule
\end{tabularx}

\subsubsection{Seuils (par véhicule)}
Stockés séparément pour Thaiger 7 et Bengalo :

\begin{tabularx}{\linewidth}{@{}llL@{}}
\toprule
\textbf{Réglage} & \textbf{Plage} & \textbf{Défaut (Thaiger 7 / Bengalo)}\\
\midrule
Température pile max & 30--120 \degC{} & 70 \degC{} / 65 \degC{}\\
Écart de tension de cellule max & 10--500 mV & 50 mV / 50 mV\\
\bottomrule
\end{tabularx}

\subsection{Écran 6 --- Réglages MQTT}
Un sous-écran des réglages qui configure le relais vers le portable ingénieur. Le
relais ne tourne que s'il est activé \emph{avant} le démarrage d'une session.

\begin{tabularx}{\linewidth}{@{}llL@{}}
\toprule
\textbf{Champ} & \textbf{Défaut} & \textbf{Remarques}\\
\midrule
Activer le relais & Désactivé & Interrupteur principal --- à activer pour publier\\
Utiliser TLS & Activé & Connexion chiffrée (exigée par HiveMQ Cloud)\\
Hôte du courtier & --- & p.\,ex. \code{abc123.s1.eu.hivemq.cloud}\\
Port & 8883 & Port MQTT TLS (8883 typique ; 1883 en clair)\\
Nom d'utilisateur & --- & Identifiant du courtier\\
Mot de passe & --- & Mot de passe (affiché masqué)\\
Fréquence de publication & 5 Hz & Cadence de publication, 1--10 Hz\\
\bottomrule
\end{tabularx}
\begin{tgnote}
Le téléphone publie en \textbf{MQTT natif (TLS)}. Le tableau de bord ingénieur se
connecte au \emph{même} courtier en \textbf{MQTT-sur-WebSocket} --- le courtier doit
donc aussi exposer un écouteur WebSocket (HiveMQ Cloud : port 8884).
\end{tgnote}

\section{Le tableau de bord ingénieur}
\code{engineer\_dashboard.html} est un fichier unique et autonome. Aucune
construction ni serveur --- juste un navigateur moderne et un accès internet (tuiles
de carte, polices et bibliothèques MQTT/graphiques via CDN).

\subsection{Ouverture \& connexion}
\begin{enumerate}
  \item Ouvrez \code{engineer\_dashboard.html} dans Chrome/Edge/Firefox.
  \item La boîte \textbf{Connection Settings} s'ouvre automatiquement la première
    fois (ou cliquez sur l'engrenage, en haut à droite).
  \item Renseignez : \textbf{Car} (\code{Thaiger 7} ou \code{Bengalo}, identique au
    véhicule publié) ; \textbf{Broker URL (WebSocket)}, p.\,ex.
    \code{wss://abc123.s1.eu.hivemq.cloud:\allowbreak8884/mqtt} (\code{wss://} pour TLS, ou
    \code{ws://host:port} en local) ; \textbf{Username} / \textbf{Password}.
  \item Cliquez \textbf{Connect}. Les réglages sont mémorisés dans le navigateur
    (\code{localStorage}) ; reconnexion automatique ensuite.
\end{enumerate}
Le point d'état (en haut à droite) : gris déconnecté, ambre en cours, vert
connecté. Trois onglets basculent entre \textbf{Operations}, \textbf{Live Tracking}
et \textbf{Analytics}.

\subsection{Vue 1 --- Operations}
Un aperçu d'un coup d'œil : une \textbf{carte vitesse principale} (bleue à/au-dessus
de la vitesse optimale, rouge en dessous) avec vitesse moyenne, nombre de tours,
temps de session, tour cible et distance ; une grille \textbf{Électrique} ; une
grille \textbf{Thermique \& Rendement} ; et une pastille \textbf{LIVE} clignotante.

\subsection{Vue 2 --- Live Tracking \& carte}
Une \textbf{OpenStreetMap} sombre plein écran (Leaflet) qui suit la voiture.
\begin{itemize}
  \item \textbf{Trace colorée par la puissance} --- dessinée segment par segment et
    colorée selon la puissance moteur : \textcolor{tgreen}{vert (0 W)} \arrow{}
    \textcolor{tamber}{ambre (100 W)} \arrow{} \textcolor{tred}{rouge (200 W+)}. La
    partie la plus récente est opaque puis s'estompe. (200 W est la puissance max de
    suivi ThaiGer ; la constante \code{POWER\_MAX} est modifiable.)
  \item \textbf{Légende de puissance} (en bas à gauche) --- l'échelle de couleurs
    plus un affichage en watts en direct.
  \item \textbf{Marqueur de voiture} orienté dans le sens de la marche.
  \item \textbf{Panneau d'état GPS} (à droite) --- qualité du fix, précision,
    latitude, longitude, altitude, vitesse GPS, cap, puissance moteur, nombre de
    tours et un compteur \textbf{Filtered}.
\end{itemize}

\subsubsection{Nettoyage du bruit GPS}
Chaque fix est nettoyé en quatre étapes : (1) rejet des fix de précision pire que
75 m ; (2) rejet des « téléportations » --- vitesse implicite impossible (plus de
162 km/h env.), avec une garde de resynchronisation ; (3) ignorer le tremblement
sous 1,5 m à l'arrêt ; (4) léger lissage. Résultat : une trace propre sans sauts.

\subsubsection{Capture par tour + remise à zéro}
À chaque incrément du compteur de tours, le tableau de bord rend la trace du tour
terminé en image PNG et la télécharge
(\code{thaiger-lapN-<horodatage>.png}), puis efface la trace. Un message vert
confirme chaque enregistrement.

\subsection{Vue 3 --- Analytics}
Un \textbf{tableau de télémétrie en direct} listant les 22 champs, plus un graphique
\textbf{Speed} (km/h dans le temps) et un graphique \textbf{Efficiency} (rendement
pile et système), chacun sur les 60 derniers points.

\section{Protocole de télémétrie \& référence des champs}
Le calculateur envoie des trames ASCII par BLE. Chaque champ est encadré
d'astérisques avec une clé d'une lettre :

\begin{lstlisting}
*A28.3**B25.0**C3**D15:35**N53.4*
\end{lstlisting}
L'analyseur tolère le flux et conserve toujours la dernière valeur connue de chaque
champ.

\begin{tabularx}{\linewidth}{@{}clL@{}}
\toprule
\textbf{Clé} & \textbf{Champ} & \textbf{Unité}\\
\midrule
A & Vitesse & km/h\\
B & Vitesse moyenne & km/h\\
C & Nombre de tours & ---\\
D & Temps total & mm:ss\\
E & Temps de tour cible & mm:ss\\
F & Vitesse optimale & km/h\\
G & Tension pile & V\\
H & Tension supercondensateur & V\\
I & Tension moteur & V\\
J & Courant pile & A\\
K & Courant supercondensateur & A\\
L & Courant moteur & A\\
M & Consommation propre & A\\
N & Température pile & \degC{}\\
O & Rapport cyclique pompe à air & \%\\
P & Conseil de conduite & ---\\
Q & Écart de tension de cellule & mV\\
R & Énergie pile (cumulée) & Ws\\
S & Énergie moteur (cumulée) & Ws\\
T & Rendement pile & \%\\
U & Rendement système & \%\\
\bottomrule
\end{tabularx}

\medskip
La distance, la puissance moteur ($V \times A$) et l'énergie en Wh sont calculées sur
le téléphone.

\section{Sujets MQTT \& charges utiles}
Tous les messages sont en \textbf{QoS 0}. Les champs vides/NaN sont omis (charge
utile d'environ 100--200 octets).

\noindent\textbf{\code{thaiger/\{car\}/telemetry}} --- publié à la cadence réglée
(5 Hz par défaut) :

\begin{lstlisting}
{"ts":1736517201234,"A":28.3,"B":25.0,"C":3,"D":"15:35","G":28.3,"N":53.4,"dist":0.42}
\end{lstlisting}
\noindent\textbf{\code{thaiger/\{car\}/gps}} --- publié quand un fix GPS est
disponible :

\begin{lstlisting}
{"ts":1736517201300,"lat":13.756331,"lon":100.501762,"alt":12.0,"spd":53.4,"acc":3.5,"hdg":270.0}
\end{lstlisting}

\begin{tabularx}{\linewidth}{@{}lLl@{}}
\toprule
\textbf{Champ} & \textbf{Signification} & \textbf{Unité}\\
\midrule
ts & Horodatage Unix & ms\\
lat / lon & Latitude / longitude & deg\\
alt & Altitude (si disponible) & m\\
spd & Vitesse GPS (si disponible) & km/h\\
acc & Précision horizontale (si disponible) & m\\
hdg & Cap (si disponible) & deg\\
\bottomrule
\end{tabularx}

\section{Profils de véhicule}

\begin{tabularx}{\linewidth}{@{}llllL@{}}
\toprule
\textbf{Profil} & \textbf{id} & \textbf{Puiss. max} & \textbf{Limite temp. pile} & \textbf{Vitesse cible min}\\
\midrule
Thaiger 7 & \code{thaiger7} & 600 W & 70 \degC{} & 25 km/h\\
Bengalo & \code{bengalo} & 900 W & 65 \degC{} & 25 km/h\\
\bottomrule
\end{tabularx}

\medskip
Les seuils sont réglables par voiture dans les \textbf{Réglages}. La trace de la
carte ingénieur utilise une échelle de puissance distincte de \textbf{200 W}
(\code{POWER\_MAX}).

\section{Dépannage}

\begin{tabularx}{\linewidth}{@{}LL@{}}
\toprule
\textbf{Symptôme} & \textbf{Cause / solution}\\
\midrule
Pas de connexion Bluetooth & Appairez d'abord le module dans les réglages Bluetooth ; vérifiez l'alimentation ; accordez Bluetooth + Localisation. Utilisez \textbf{Pair New}.\\
« Connected » mais pas de données & Mauvais appareil, ou calculateur muet. Testez le mode démo, puis revérifiez le module.\\
Le tableau de bord ne passe pas au vert & Mauvaise URL WebSocket (doit être \code{wss://\ldots:8884/mqtt}, pas le port \code{8883}), identifiants faux, ou pas d'écouteur WebSocket.\\
Vert mais pas de télémétrie & Le \textbf{Car} du tableau de bord ne correspond pas à l'id du téléphone, ou le relais n'est pas activé. Même \code{thaiger7}/\code{bengalo} des deux côtés.\\
Carte vide & Pas d'internet pour les tuiles OSM, ou pas encore de fix GPS. La carte se centre au premier fix.\\
La trace GPS saute & Le nettoyage devrait l'éviter ; pour des fix très grossiers, augmentez \code{GPS\_MAX\_ACC}. Le compteur \textbf{Filtered} indique les rejets.\\
GPS lent (quelques secondes) & Vue dégagée du ciel et autorisation \textbf{Localisation} ; en intérieur, repli sur des fix réseau grossiers.\\
Pas de CSV & Terminez une session avec des trames enregistrées ; sinon \emph{Export CSV} indique « No data to export ».\\
\bottomrule
\end{tabularx}

\section{Réglages avancés}
Sur le téléphone (Réglages / Réglages MQTT) : écran allumé, couleur de vitesse,
cadence de l'IU, vibreur/son, seuils par voiture, et tous les paramètres
courtier/relais.

\medskip
\noindent Dans \code{engineer\_dashboard.html} (constantes en haut du script) :

\begin{tabularx}{\linewidth}{@{}llL@{}}
\toprule
\textbf{Constante} & \textbf{Défaut} & \textbf{Rôle}\\
\midrule
\code{POWER\_MAX} & 200 & Valeur en watts mappée au rouge plein\\
\code{GPS\_MAX\_ACC} & 75 & Rejeter les fix moins précis que ce nombre de mètres\\
\code{GPS\_MAX\_SPEED} & 45 & m/s ; au-delà, considéré comme téléportation\\
\code{GPS\_MIN\_MOVE} & 1.5 & m ; en dessous, tremblement à l'arrêt\\
\code{GPS\_SMOOTH} & 0.35 & Facteur de lissage EMA (0 = aucun, 1 = brut)\\
\code{GPS\_MAX\_REJECTS} & 4 & Resync après ce nombre de rejets consécutifs\\
\code{TRAIL\_MAX} & 250 & Nombre max de points de trace en mémoire\\
\bottomrule
\end{tabularx}

\section{Glossaire}
\begin{description}[style=nextline,leftmargin=0pt,labelwidth=0pt,labelsep=0pt,itemindent=0pt]
  \item[BLE] Bluetooth Low Energy, la liaison radio du calculateur au téléphone.
  \item[MQTT] Protocole léger de publication/abonnement ; le courtier relaie les messages.
  \item[Courtier (broker)] Le serveur MQTT (p.\,ex. HiveMQ Cloud) auquel les deux côtés se connectent.
  \item[TLS] Chiffrement de la connexion au courtier.
  \item[WebSocket] Le transport utilisé par le navigateur pour parler MQTT (\code{ws://} / \code{wss://}).
  \item[Pile à combustible (FC)] La source d'énergie à hydrogène de la voiture.
  \item[Supercondensateur] Tampon d'énergie pour les pointes de charge.
  \item[Trace] La ligne dessinée derrière la voiture, colorée par la puissance moteur.
\end{description}


% #####################################################################
%  PART V - ESPAÑOL
% #####################################################################
\selectlanguage{spanish}
\guidepart{Guía en español}

\section{Qué es este proyecto}
ThaiGer Track Control es un sistema de telemetría de dos partes para un coche de
carreras de pila de combustible de hidrógeno:
\begin{itemize}
  \item Una \textbf{app de Android} que se ejecuta en un teléfono \emph{dentro del
    coche}. Recibe la telemetría en vivo del controlador del coche por
    \textbf{Bluetooth Low Energy (BLE)}, la muestra en un panel de carreras para el
    piloto, registra la sesión y \textbf{retransmite} los datos por internet vía
    \textbf{MQTT}.
  \item Un \textbf{panel del ingeniero} --- un único archivo HTML
    (\code{engineer\_dashboard.html}) --- que se ejecuta en un navegador en el
    \emph{portátil del ingeniero} en el box. Se suscribe al mismo flujo MQTT y
    muestra datos de operación, un mapa en vivo con una traza GPS coloreada por
    potencia, y gráficas de análisis.
\end{itemize}
Las dos mitades nunca se comunican directamente; se encuentran en el
\textbf{broker MQTT} en la nube.

\section{Visión general del sistema}

\begin{lstlisting}
   +------------+   BLE    +------------------+   MQTT/TLS   +----------------+
   | Car        | -------> | Android app      | -----------> | MQTT broker    |
   | controller | *A28.3*  | (phone in car)   | JSON frames  | (HiveMQ Cloud) |
   +------------+          +------------------+              +-------+--------+
                                  |                                  | WebSocket/TLS
                                  | GPS (phone LocationManager)      v
                                  v                          +----------------+
                          thaiger/{car}/gps                  | Engineer       |
                          thaiger/{car}/telemetry            | Dashboard HTML |
                                                             +----------------+
\end{lstlisting}
\begin{itemize}
  \item Dos temas MQTT por coche: \code{thaiger/\{car\}/telemetry} y
    \code{thaiger/\{car\}/gps}.
  \item \code{\{car\}} es el id del vehículo: \code{thaiger7} o \code{bengalo}.
  \item El teléfono publica; el panel del ingeniero se suscribe. Cualquier número de
    portátiles puede observar el mismo coche a la vez.
\end{itemize}

\section{La app de Android}

\subsection{Instalación \& primer arranque}
La app es un proyecto Android estándar (Gradle). Para construir un APK de
depuración:

\begin{lstlisting}
./gradlew assembleDebug
\end{lstlisting}

\begin{tabularx}{\linewidth}{@{}lL@{}}
\toprule
\textbf{Elemento} & \textbf{Valor}\\
\midrule
Android mínimo & 9.0 (API 28)\\
Android objetivo & 14 (API 34)\\
Lenguaje & Java 11\\
Id de app & \code{com.thaiger.h2racing}\\
\bottomrule
\end{tabularx}

\medskip
Instala el APK en el teléfono que va en el coche. En el primer arranque, la app
abre la pantalla \textbf{Selección de vehículo}.
\begin{tgnote}
\textbf{Consejo.} Monta el teléfono en horizontal, visible para el piloto. El panel
mantiene la pantalla encendida durante una sesión (Ajustes \arrow{} \emph{Mantener
pantalla encendida}).
\end{tgnote}

\subsection{Permisos}
La app pide permisos en la pantalla \textbf{Conexión}, antes de usar el hardware.
Concédelos todos para una funcionalidad completa.

\begin{tabularx}{\linewidth}{@{}lL@{}}
\toprule
\textbf{Permiso} & \textbf{Por qué es necesario}\\
\midrule
Bluetooth Connect / Scan (Android 12+) & Emparejar \& hablar con el módulo BLE\\
Bluetooth / Bluetooth Admin (Android $\le$ 11) & Igual, en Android más antiguo\\
Ubicación (precisa) & Exigida por Android para el escaneo BLE \emph{y} el GPS\\
Ubicación (aproximada) & Respaldo GPS\\
Internet & Relé MQTT al portátil del ingeniero\\
Wake Lock & Mantener la pantalla encendida durante una sesión\\
Vibración & Respuesta háptica para alertas\\
\bottomrule
\end{tabularx}

\medskip
Si rechazas el Bluetooth, aún puedes usar el \textbf{modo demo} (datos sintéticos).

\subsection{Pantalla 1 --- Selección de vehículo}
\begin{itemize}
  \item Dos tarjetas: \textbf{Thaiger 7} y \textbf{Bengalo}. Toca una tarjeta para
    seleccionarla (un borde azul marca la selección). Thaiger 7 por defecto.
  \item El perfil elegido controla los umbrales de alerta y la escala de potencia
    máxima.
  \item Toca \textbf{Settings} (arriba) para abrir los ajustes antes de una sesión.
  \item Toca \textbf{Connect to \ldots} para pasar a la pantalla Conexión.
\end{itemize}

\subsection{Pantalla 2 --- Conexión (Bluetooth / modo demo)}
Esta pantalla empareja el teléfono con el coche e inicia el flujo de datos.
\begin{enumerate}
  \item La app comprueba los permisos y que el Bluetooth esté activado.
  \item Un diálogo lista tus \textbf{dispositivos Bluetooth emparejados} más una
    entrada final \textbf{modo demo}.
  \item Elige el módulo de tu coche (p.\,ej. un dispositivo llamado HM-10 / HC-05 /
    ESP32), \textbf{o} elige el \textbf{modo demo} para un flujo sintético sin
    hardware.
  \item \textbf{Pair New} abre los ajustes Bluetooth de Android si el módulo no está
    emparejado. \textbf{Cancel} vuelve a la selección de vehículo.
\end{enumerate}
Cuando el módulo indica \textbf{Connected}, la app avanza automáticamente al panel
en vivo. Si el relé MQTT está activado, también arranca aquí: el portátil del
ingeniero recibe datos en cuanto empieza la sesión.
\begin{tgnote}
Las caídas de conexión se gestionan automáticamente con retroceso exponencial
(1 \arrow{} 2 \arrow{} 4 \arrow{} 8 \arrow{} 16 s). Verás \emph{Reconnecting\ldots}
--- sin acción necesaria.
\end{tgnote}

\subsection{Pantalla 3 --- Panel en vivo}
El panel es la vista principal del piloto (horizontal, pantalla encendida). Muestra
en continuo el último valor de cada campo de telemetría.
\begin{itemize}
  \item \textbf{Velocidad} --- el número principal grande (km/h). Si la \emph{alerta
    de color de velocidad} está activa, el número se vuelve \textbf{rojo} por debajo
    de la velocidad objetivo mínima. Un indicador muestra \ok{} on target /
    \nope{} too slow.
  \item \textbf{Barra de vueltas} (una línea sobre la barra inferior):
    \code{PREV 1:18} \midsep{} \code{LAP 3} \midsep{} \code{CURR 0:09} --- tiempo de
    la vuelta anterior, número de vuelta actual y tiempo actual.
  \item \textbf{Bloque eléctrico} --- tensión \& corriente de pila /
    supercondensador / motor, consumo propio, diferencia de tensión de celda.
  \item \textbf{Térmica \& eficiencia} --- temperatura de pila, ciclo de trabajo de
    la bomba de aire, eficiencia de pila y sistema, energías, velocidad óptima,
    consejo de conducción.
  \item \textbf{Banner de alerta} --- aparece al cruzar un umbral (p.\,ej.
    temperatura de pila o diferencia de tensión demasiado altas). Con
    \emph{Vibración} / \emph{Sonido de alerta}, respuesta háptica / sonora además.
    Toca \textbf{Dismiss} para ocultarlo.
\end{itemize}
\begin{tgnote}
\textbf{Terminar una sesión:} \textbf{pulsación larga sobre el número de
velocidad.} Detiene el Bluetooth y el relé MQTT, finaliza las estadísticas y abre el
\textbf{resumen de sesión}.
\end{tgnote}

\subsection{Pantalla 4 --- Resumen de sesión \& exportación CSV}
Tras una sesión verás resultados agregados: duración, distancia, velocidad media,
energía consumida, potencia máxima del motor, temperatura máxima de pila y número de
alertas, más una \textbf{gráfica de potencia}.
\begin{itemize}
  \item \textbf{Export CSV} --- escribe cada trama registrada en un archivo CSV y
    abre el panel de compartir de Android. El archivo se llama
    \code{thaiger\_\{car\}\_\allowbreak\{YYYYMMDD\_HHMM\}.csv} y tiene 21 columnas.
    Sin datos aparece \emph{«No data to export»}.
  \item \textbf{New Run} --- vuelta directa a la Conexión.
  \item \textbf{Back to Car Select} --- volver al inicio.
\end{itemize}

\subsection{Pantalla 5 --- Ajustes}
Accesible desde \textbf{Settings} en la selección de vehículo. Todo se guarda al
instante.

\subsubsection{Pantalla}

\begin{tabularx}{\linewidth}{@{}llL@{}}
\toprule
\textbf{Ajuste} & \textbf{Defecto} & \textbf{Significado}\\
\midrule
Mantener pantalla encendida & Sí & Mantiene un wake lock para que la pantalla no se apague\\
Alerta de color de velocidad & Sí & Velocidad en rojo bajo la velocidad objetivo mínima\\
Frecuencia de actualización & 50 ms & Cada cuánto se redibuja la IU (20--500 ms)\\
\bottomrule
\end{tabularx}

\subsubsection{Alertas}

\begin{tabularx}{\linewidth}{@{}llL@{}}
\toprule
\textbf{Ajuste} & \textbf{Defecto} & \textbf{Significado}\\
\midrule
Vibración & Sí & Impulso háptico en cada alerta\\
Sonido de alerta & No & Reproduce un tono en cada alerta\\
\bottomrule
\end{tabularx}

\subsubsection{Umbrales (por vehículo)}
Guardados por separado para Thaiger 7 y Bengalo:

\begin{tabularx}{\linewidth}{@{}llL@{}}
\toprule
\textbf{Ajuste} & \textbf{Rango} & \textbf{Defecto (Thaiger 7 / Bengalo)}\\
\midrule
Temperatura pila máx & 30--120 \degC{} & 70 \degC{} / 65 \degC{}\\
Diferencia de tensión de celda máx & 10--500 mV & 50 mV / 50 mV\\
\bottomrule
\end{tabularx}

\subsection{Pantalla 6 --- Ajustes MQTT}
Una subpantalla de ajustes que configura el relé al portátil del ingeniero. El relé
solo funciona si se activa \emph{antes} de iniciar una sesión.

\begin{tabularx}{\linewidth}{@{}llL@{}}
\toprule
\textbf{Campo} & \textbf{Defecto} & \textbf{Notas}\\
\midrule
Activar relé & No & Interruptor principal --- actívalo para publicar\\
Usar TLS & Sí & Conexión cifrada (exigida por HiveMQ Cloud)\\
Host del broker & --- & p.\,ej. \code{abc123.s1.eu.hivemq.cloud}\\
Puerto & 8883 & Puerto MQTT TLS (8883 típico; 1883 en claro)\\
Usuario & --- & Usuario del broker\\
Contraseña & --- & Contraseña (mostrada oculta)\\
Frecuencia de publicación & 5 Hz & Tasa de publicación, 1--10 Hz\\
\bottomrule
\end{tabularx}
\begin{tgnote}
El teléfono publica por \textbf{MQTT nativo (TLS)}. El panel del ingeniero se conecta
al \emph{mismo} broker por \textbf{MQTT sobre WebSocket} --- el broker debe exponer
también un escuchador WebSocket (HiveMQ Cloud: puerto 8884).
\end{tgnote}

\section{El panel del ingeniero}
\code{engineer\_dashboard.html} es un único archivo autónomo. Sin compilación ni
servidor --- solo un navegador moderno y acceso a internet (teselas del mapa, fuentes
y bibliotecas MQTT/gráficas vía CDN).

\subsection{Apertura \& conexión}
\begin{enumerate}
  \item Abre \code{engineer\_dashboard.html} en Chrome/Edge/Firefox.
  \item El diálogo \textbf{Connection Settings} se abre solo la primera vez (o haz
    clic en el engranaje, arriba a la derecha).
  \item Rellena: \textbf{Car} (\code{Thaiger 7} o \code{Bengalo}, igual al vehículo
    publicado); \textbf{Broker URL (WebSocket)}, p.\,ej.
    \code{wss://abc123.s1.eu.hivemq.cloud:\allowbreak8884/mqtt} (\code{wss://} para TLS, o
    \code{ws://host:port} en local); \textbf{Username} / \textbf{Password}.
  \item Haz clic en \textbf{Connect}. Los ajustes se guardan en el navegador
    (\code{localStorage}); reconexión automática la próxima vez.
\end{enumerate}
El punto de estado (arriba a la derecha): gris desconectado, ámbar conectando, verde
conectado. Tres pestañas cambian entre \textbf{Operations}, \textbf{Live Tracking} y
\textbf{Analytics}.

\subsection{Vista 1 --- Operations}
Una visión de un vistazo: una \textbf{tarjeta de velocidad principal} (azul a/por
encima de la velocidad óptima, roja por debajo) con velocidad media, vueltas, tiempo
de sesión, vuelta objetivo y distancia; una rejilla \textbf{Eléctrica}; una rejilla
\textbf{Térmica \& Eficiencia}; y una píldora \textbf{LIVE} parpadeante.

\subsection{Vista 2 --- Live Tracking \& mapa}
Un \textbf{OpenStreetMap} oscuro a pantalla completa (Leaflet) que sigue al coche.
\begin{itemize}
  \item \textbf{Traza coloreada por potencia} --- dibujada segmento a segmento y
    coloreada según la potencia del motor: \textcolor{tgreen}{verde (0 W)} \arrow{}
    \textcolor{tamber}{ámbar (100 W)} \arrow{} \textcolor{tred}{rojo (200 W+)}. La
    parte más reciente es opaca y se desvanece con el tiempo. (200 W es la potencia
    máxima de seguimiento ThaiGer; la constante \code{POWER\_MAX} es modificable.)
  \item \textbf{Leyenda de potencia} (abajo a la izquierda) --- la escala de color
    más una lectura en vatios en vivo.
  \item \textbf{Marcador del coche} orientado en el sentido de la marcha.
  \item \textbf{Panel de estado GPS} (derecha) --- calidad del fix, precisión,
    latitud, longitud, altitud, velocidad GPS, rumbo, potencia del motor, vueltas y
    un contador \textbf{Filtered}.
\end{itemize}

\subsubsection{Limpieza del ruido GPS}
Cada fix se limpia en cuatro pasos: (1) descartar fixes con precisión peor que 75 m;
(2) rechazar «teletransportes» --- velocidad implícita imposible (más de 162 km/h
aprox.), con una guarda de resincronización; (3) ignorar el temblor bajo 1,5 m en
parado; (4) suavizado ligero. Resultado: una traza limpia sin saltos.

\subsubsection{Captura por vuelta + reinicio}
Cada vez que el contador de vueltas se incrementa, el panel renderiza la traza de la
vuelta terminada como imagen PNG y la descarga (\code{thaiger-lapN-<marca>.png}),
y luego borra la traza. Un aviso verde confirma cada guardado.

\subsection{Vista 3 --- Analytics}
Una \textbf{tabla de telemetría en vivo} con los 22 campos, más una gráfica
\textbf{Speed} (km/h en el tiempo) y una gráfica \textbf{Efficiency} (eficiencia de
pila y sistema), cada una con los últimos 60 puntos.

\section{Protocolo de telemetría \& referencia de campos}
El controlador envía tramas ASCII por BLE. Cada campo va entre asteriscos con una
clave de una letra:

\begin{lstlisting}
*A28.3**B25.0**C3**D15:35**N53.4*
\end{lstlisting}
El analizador tolera el flujo y conserva siempre el último valor conocido de cada
campo.

\begin{tabularx}{\linewidth}{@{}clL@{}}
\toprule
\textbf{Clave} & \textbf{Campo} & \textbf{Unidad}\\
\midrule
A & Velocidad & km/h\\
B & Velocidad media & km/h\\
C & Número de vueltas & ---\\
D & Tiempo total & mm:ss\\
E & Tiempo de vuelta objetivo & mm:ss\\
F & Velocidad óptima & km/h\\
G & Tensión de pila & V\\
H & Tensión de supercondensador & V\\
I & Tensión de motor & V\\
J & Corriente de pila & A\\
K & Corriente de supercondensador & A\\
L & Corriente de motor & A\\
M & Consumo propio & A\\
N & Temperatura de pila & \degC{}\\
O & Ciclo de trabajo de bomba de aire & \%\\
P & Consejo de conducción & ---\\
Q & Diferencia de tensión de celda & mV\\
R & Energía de pila (acumulada) & Ws\\
S & Energía de motor (acumulada) & Ws\\
T & Eficiencia de pila & \%\\
U & Eficiencia del sistema & \%\\
\bottomrule
\end{tabularx}

\medskip
La distancia, la potencia del motor ($V \times A$) y la energía en Wh se calculan en
el teléfono.

\section{Temas MQTT \& cargas útiles}
Todos los mensajes son \textbf{QoS 0}. Los campos vacíos/NaN se omiten (carga útil de
unos 100--200 bytes).

\noindent\textbf{\code{thaiger/\{car\}/telemetry}} --- publicado a la tasa
configurada (5 Hz por defecto):

\begin{lstlisting}
{"ts":1736517201234,"A":28.3,"B":25.0,"C":3,"D":"15:35","G":28.3,"N":53.4,"dist":0.42}
\end{lstlisting}
\noindent\textbf{\code{thaiger/\{car\}/gps}} --- publicado cuando hay un fix GPS:

\begin{lstlisting}
{"ts":1736517201300,"lat":13.756331,"lon":100.501762,"alt":12.0,"spd":53.4,"acc":3.5,"hdg":270.0}
\end{lstlisting}

\begin{tabularx}{\linewidth}{@{}lLl@{}}
\toprule
\textbf{Campo} & \textbf{Significado} & \textbf{Unidad}\\
\midrule
ts & Marca de tiempo Unix & ms\\
lat / lon & Latitud / longitud & grados\\
alt & Altitud (si disponible) & m\\
spd & Velocidad GPS (si disponible) & km/h\\
acc & Precisión horizontal (si disponible) & m\\
hdg & Rumbo (si disponible) & grados\\
\bottomrule
\end{tabularx}

\section{Perfiles de vehículo}

\begin{tabularx}{\linewidth}{@{}llllL@{}}
\toprule
\textbf{Perfil} & \textbf{id} & \textbf{Pot. máx} & \textbf{Límite temp. pila} & \textbf{Vel. objetivo mín}\\
\midrule
Thaiger 7 & \code{thaiger7} & 600 W & 70 \degC{} & 25 km/h\\
Bengalo & \code{bengalo} & 900 W & 65 \degC{} & 25 km/h\\
\bottomrule
\end{tabularx}

\medskip
Los umbrales se pueden ajustar por coche en los \textbf{Ajustes}. La traza del mapa
del ingeniero usa una escala de potencia distinta de \textbf{200 W}
(\code{POWER\_MAX}).

\section{Resolución de problemas}

\begin{tabularx}{\linewidth}{@{}LL@{}}
\toprule
\textbf{Síntoma} & \textbf{Causa / solución}\\
\midrule
No conecta por Bluetooth & Empareja primero el módulo en los ajustes Bluetooth; comprueba la alimentación; concede Bluetooth + Ubicación. Usa \textbf{Pair New}.\\
«Connected» pero sin datos & Dispositivo equivocado, o el controlador no envía. Prueba el modo demo y revisa el módulo.\\
El panel nunca se pone verde & URL WebSocket incorrecta (debe ser \code{wss://\ldots:8884/mqtt}, no el puerto \code{8883}), credenciales erróneas, o sin escuchador WebSocket.\\
Verde pero sin telemetría & El \textbf{Car} del panel no coincide con el id del teléfono, o el relé no está activado. Mismo \code{thaiger7}/\code{bengalo} en ambos.\\
Mapa en blanco & Sin internet para las teselas OSM, o sin fix GPS aún. El mapa se centra al primer fix.\\
La traza GPS salta & La limpieza debería evitarlo; con fixes muy gruesos, sube \code{GPS\_MAX\_ACC}. El contador \textbf{Filtered} muestra los descartes.\\
GPS lento (cada varios segundos) & Asegura vista despejada del cielo y permiso de \textbf{Ubicación}; en interior recurre a fixes de red gruesos.\\
No se genera CSV & Completa una sesión con tramas registradas; si no, \emph{Export CSV} indica «No data to export».\\
\bottomrule
\end{tabularx}

\section{Ajustes avanzados}
En el teléfono (Ajustes / Ajustes MQTT): pantalla encendida, color de velocidad,
tasa de la IU, vibración/sonido, umbrales por coche, y todos los parámetros de
broker/relé.

\medskip
\noindent En \code{engineer\_dashboard.html} (constantes al inicio del script):

\begin{tabularx}{\linewidth}{@{}llL@{}}
\toprule
\textbf{Constante} & \textbf{Defecto} & \textbf{Propósito}\\
\midrule
\code{POWER\_MAX} & 200 & Valor en vatios mapeado al rojo pleno\\
\code{GPS\_MAX\_ACC} & 75 & Descartar fixes peores que estos metros\\
\code{GPS\_MAX\_SPEED} & 45 & m/s; más rápido se trata como teletransporte\\
\code{GPS\_MIN\_MOVE} & 1.5 & m; movimientos menores son temblor en parado\\
\code{GPS\_SMOOTH} & 0.35 & Factor de suavizado EMA (0 = ninguno, 1 = crudo)\\
\code{GPS\_MAX\_REJECTS} & 4 & Resync tras tantos rechazos seguidos\\
\code{TRAIL\_MAX} & 250 & Máx. de puntos de traza en memoria\\
\bottomrule
\end{tabularx}

\section{Glosario}
\begin{description}[style=nextline,leftmargin=0pt,labelwidth=0pt,labelsep=0pt,itemindent=0pt]
  \item[BLE] Bluetooth Low Energy, el enlace de radio del controlador al teléfono.
  \item[MQTT] Protocolo ligero de publicación/suscripción; el broker reenvía los mensajes.
  \item[Broker] El servidor MQTT (p.\,ej. HiveMQ Cloud) al que se conectan ambos lados.
  \item[TLS] Cifrado de la conexión al broker.
  \item[WebSocket] El transporte que usa el navegador para hablar MQTT (\code{ws://} / \code{wss://}).
  \item[Pila de combustible (FC)] La fuente de energía de hidrógeno del coche.
  \item[Supercondensador] Búfer de energía para los picos de carga.
  \item[Traza] La línea dibujada tras el coche, coloreada por la potencia del motor.
\end{description}


% #####################################################################
%  PART VI - РУССКИЙ
% #####################################################################
\selectlanguage{russian}
\guidepart{Руководство пользователя}

\section{Что это за проект}
ThaiGer Track Control --- это система телеметрии из двух частей для гоночного
автомобиля на водородных топливных элементах:
\begin{itemize}
  \item \textbf{Android-приложение}, работающее на телефоне \emph{внутри
    автомобиля}. Оно получает телеметрию в реальном времени от контроллера машины
    по \textbf{Bluetooth Low Energy (BLE)}, показывает её на гоночной панели для
    пилота, записывает заезд и \textbf{ретранслирует} данные через интернет по
    \textbf{MQTT}.
  \item \textbf{Панель инженера} --- один HTML-файл
    (\code{engineer\_dashboard.html}) --- работающая в браузере на \emph{ноутбуке
    инженера} в боксе. Она подписывается на тот же поток MQTT и показывает
    эксплуатационные данные, живую карту с GPS-следом, окрашенным по мощности, и
    графики анализа.
\end{itemize}
Две половины никогда не общаются напрямую; они встречаются на \textbf{брокере MQTT}
в облаке.

\section{Обзор системы}

\begin{lstlisting}
   +------------+   BLE    +------------------+   MQTT/TLS   +----------------+
   | Car        | -------> | Android app      | -----------> | MQTT broker    |
   | controller | *A28.3*  | (phone in car)   | JSON frames  | (HiveMQ Cloud) |
   +------------+          +------------------+              +-------+--------+
                                  |                                  | WebSocket/TLS
                                  | GPS (phone LocationManager)      v
                                  v                          +----------------+
                          thaiger/{car}/gps                  | Engineer       |
                          thaiger/{car}/telemetry            | Dashboard HTML |
                                                             +----------------+
\end{lstlisting}
\begin{itemize}
  \item Два топика MQTT на каждую машину: \code{thaiger/\{car\}/telemetry} и
    \code{thaiger/\{car\}/gps}.
  \item \code{\{car\}} --- идентификатор автомобиля: \code{thaiger7} или
    \code{bengalo}.
  \item Телефон публикует; панель инженера подписывается. Любое число ноутбуков
    может одновременно наблюдать за одной машиной.
\end{itemize}

\section{Android-приложение}

\subsection{Установка \& первый запуск}
Приложение --- стандартный проект Android (Gradle). Сборка отладочного APK:

\begin{lstlisting}
./gradlew assembleDebug
\end{lstlisting}

\begin{tabularx}{\linewidth}{@{}lL@{}}
\toprule
\textbf{Пункт} & \textbf{Значение}\\
\midrule
Минимум Android & 9.0 (API 28)\\
Целевой Android & 14 (API 34)\\
Язык & Java 11\\
Идентификатор & \code{com.thaiger.h2racing}\\
\bottomrule
\end{tabularx}

\medskip
Установите APK на телефон, который остаётся в машине. При первом запуске
открывается экран \textbf{Выбор автомобиля}.
\begin{tgnote}
\textbf{Совет.} Закрепите телефон горизонтально на виду у пилота. Панель удерживает
экран включённым во время заезда (Настройки \arrow{} \emph{Не гасить экран}).
\end{tgnote}

\subsection{Разрешения}
Приложение запрашивает разрешения на экране \textbf{Подключение}, до использования
оборудования. Предоставьте их все для полной работы.

\begin{tabularx}{\linewidth}{@{}lL@{}}
\toprule
\textbf{Разрешение} & \textbf{Зачем нужно}\\
\midrule
Bluetooth Connect / Scan (Android 12+) & Сопряжение \& связь с модулем BLE\\
Bluetooth / Bluetooth Admin (Android $\le$ 11) & То же, на старом Android\\
Местоположение (точное) & Требуется Android для сканирования BLE \emph{и} для GPS\\
Местоположение (приблизительное) & Резервный GPS\\
Интернет & Ретрансляция MQTT на ноутбук инженера\\
Wake Lock & Держать экран включённым во время заезда\\
Вибрация & Тактильная отдача для оповещений\\
\bottomrule
\end{tabularx}

\medskip
Если отказать в Bluetooth, всё равно можно использовать \textbf{демо-режим}
(синтетические данные).

\subsection{Экран 1 --- Выбор автомобиля}
\begin{itemize}
  \item Две карточки: \textbf{Thaiger 7} и \textbf{Bengalo}. Коснитесь карточки,
    чтобы выбрать (синяя рамка отмечает выбор). По умолчанию Thaiger 7.
  \item Выбранный профиль задаёт пороги предупреждений и шкалу максимальной
    мощности.
  \item Коснитесь \textbf{Settings} (вверху), чтобы открыть настройки перед заездом.
  \item Коснитесь \textbf{Connect to \ldots}, чтобы перейти к экрану Подключение.
\end{itemize}

\subsection{Экран 2 --- Подключение (Bluetooth / демо-режим)}
Этот экран связывает телефон с машиной и запускает поток данных.
\begin{enumerate}
  \item Приложение проверяет разрешения и что Bluetooth включён.
  \item Диалог перечисляет \textbf{сопряжённые устройства Bluetooth} плюс
    последний пункт \textbf{демо-режим}.
  \item Выберите модуль вашей машины (например, устройство с именем HM-10 / HC-05 /
    ESP32), \textbf{либо} выберите \textbf{демо-режим} для синтетического потока без
    оборудования.
  \item \textbf{Pair New} открывает настройки Bluetooth Android, если модуль ещё не
    сопряжён. \textbf{Cancel} возвращает к выбору автомобиля.
\end{enumerate}
Когда модуль сообщает \textbf{Connected}, приложение автоматически переходит к живой
панели. Если ретранслятор MQTT включён, он тоже стартует здесь: ноутбук инженера
начинает получать данные с началом заезда.
\begin{tgnote}
Обрывы связи обрабатываются автоматически с экспоненциальной задержкой
(1 \arrow{} 2 \arrow{} 4 \arrow{} 8 \arrow{} 16 с). Вы увидите
\emph{Reconnecting\ldots} --- действий не требуется.
\end{tgnote}

\subsection{Экран 3 --- Живая панель}
Панель --- основной вид пилота (горизонтально, экран не гаснет). Она непрерывно
показывает последнее значение каждого поля телеметрии.
\begin{itemize}
  \item \textbf{Скорость} --- большое главное число (км/ч). Если включено
    \emph{цветовое предупреждение скорости}, число становится \textbf{красным} ниже
    минимальной целевой скорости. Индикатор показывает \ok{} on target /
    \nope{} too slow.
  \item \textbf{Строка кругов} (одна строка над нижней панелью):
    \code{PREV 1:18} \midsep{} \code{LAP 3} \midsep{} \code{CURR 0:09} --- время
    предыдущего круга, номер текущего круга и текущее время.
  \item \textbf{Электрический блок} --- напряжение \& ток топливного элемента /
    суперконденсатора / двигателя, собственное потребление, разница напряжений
    ячеек.
  \item \textbf{Тепло \& эффективность} --- температура элемента, скважность
    воздушного насоса, КПД элемента и системы, энергии, оптимальная скорость,
    подсказка по вождению.
  \item \textbf{Баннер оповещения} --- появляется при превышении порога (например,
    температура элемента или разница напряжений слишком высоки). С \emph{Вибрацией}
    / \emph{Звуком оповещения} добавляется тактильная / звуковая отдача. Коснитесь
    \textbf{Dismiss}, чтобы скрыть баннер.
\end{itemize}
\begin{tgnote}
\textbf{Завершение заезда:} \textbf{долгое нажатие на число скорости.} Это
останавливает Bluetooth и ретранслятор MQTT, завершает статистику и открывает
\textbf{итог заезда}.
\end{tgnote}

\subsection{Экран 4 --- Итог заезда \& экспорт CSV}
После заезда показываются сводные результаты: длительность, дистанция, средняя
скорость, израсходованная энергия, пиковая мощность двигателя, максимальная
температура элемента и число оповещений, плюс \textbf{график мощности}.
\begin{itemize}
  \item \textbf{Export CSV} --- записывает каждый зафиксированный кадр в файл CSV и
    открывает меню «Поделиться» Android. Файл называется
    \code{thaiger\_\{car\}\_\allowbreak\{YYYYMMDD\_HHMM\}.csv} и содержит 21 столбец.
    Без данных появляется сообщение \emph{«No data to export»}.
  \item \textbf{New Run} --- сразу назад к Подключению.
  \item \textbf{Back to Car Select} --- вернуться к началу.
\end{itemize}

\subsection{Экран 5 --- Настройки}
Доступны по ссылке \textbf{Settings} на выборе автомобиля. Всё сохраняется сразу.

\subsubsection{Экран}

\begin{tabularx}{\linewidth}{@{}llL@{}}
\toprule
\textbf{Настройка} & \textbf{По умолч.} & \textbf{Значение}\\
\midrule
Не гасить экран & Вкл & Удерживает wake lock, чтобы экран не гас\\
Цветовое предупреждение скорости & Вкл & Красная скорость ниже минимальной целевой\\
Частота обновления & 50 мс & Как часто перерисовывается интерфейс (20--500 мс)\\
\bottomrule
\end{tabularx}

\subsubsection{Оповещения}

\begin{tabularx}{\linewidth}{@{}llL@{}}
\toprule
\textbf{Настройка} & \textbf{По умолч.} & \textbf{Значение}\\
\midrule
Вибрация & Вкл & Тактильный импульс при оповещении\\
Звук оповещения & Выкл & Воспроизводить тон при оповещении\\
\bottomrule
\end{tabularx}

\subsubsection{Пороги (по автомобилю)}
Хранятся отдельно для Thaiger 7 и Bengalo:

\begin{tabularx}{\linewidth}{@{}llL@{}}
\toprule
\textbf{Настройка} & \textbf{Диапазон} & \textbf{По умолч. (Thaiger 7 / Bengalo)}\\
\midrule
Макс. температура элемента & 30--120 \degC{} & 70 \degC{} / 65 \degC{}\\
Макс. разница напряжений ячеек & 10--500 мВ & 50 мВ / 50 мВ\\
\bottomrule
\end{tabularx}

\subsection{Экран 6 --- Настройки MQTT}
Подэкран настроек, конфигурирующий ретранслятор на ноутбук инженера. Ретранслятор
работает, только если включён \emph{до} начала заезда.

\begin{tabularx}{\linewidth}{@{}llL@{}}
\toprule
\textbf{Поле} & \textbf{По умолч.} & \textbf{Примечания}\\
\midrule
Включить ретранслятор & Выкл & Главный выключатель --- включите для публикации\\
Использовать TLS & Вкл & Шифрованное соединение (требуется HiveMQ Cloud)\\
Хост брокера & --- & например \code{abc123.s1.eu.hivemq.cloud}\\
Порт & 8883 & Порт MQTT TLS (обычно 8883; 1883 без шифрования)\\
Имя пользователя & --- & Имя пользователя брокера\\
Пароль & --- & Пароль брокера (показан скрытым)\\
Частота публикации & 5 Гц & Частота публикации телеметрии, 1--10 Гц\\
\bottomrule
\end{tabularx}
\begin{tgnote}
Телефон публикует по \textbf{нативному MQTT (TLS)}. Панель инженера подключается к
\emph{тому же} брокеру по \textbf{MQTT поверх WebSocket} --- поэтому брокер должен
также предоставлять слушатель WebSocket (HiveMQ Cloud: порт 8884).
\end{tgnote}

\section{Панель инженера}
\code{engineer\_dashboard.html} --- это один автономный файл. Не нужны ни сборка, ни
сервер --- только современный браузер и доступ в интернет (тайлы карты, шрифты и
библиотеки MQTT/графиков с CDN).

\subsection{Открытие \& подключение}
\begin{enumerate}
  \item Откройте \code{engineer\_dashboard.html} в Chrome/Edge/Firefox.
  \item Диалог \textbf{Connection Settings} открывается сам при первом запуске (или
    нажмите шестерёнку справа вверху).
  \item Заполните: \textbf{Car} (\code{Thaiger 7} или \code{Bengalo}, как у
    телефона); \textbf{Broker URL (WebSocket)}, например
    \code{wss://abc123.s1.eu.hivemq.cloud:\allowbreak8884/mqtt} (\code{wss://} для TLS или
    \code{ws://host:port} локально); \textbf{Username} / \textbf{Password}.
  \item Нажмите \textbf{Connect}. Настройки сохраняются в браузере
    (\code{localStorage}); в следующий раз --- автоподключение.
\end{enumerate}
Точка состояния (справа вверху): серая --- нет связи, янтарная --- подключение,
зелёная --- подключено. Три вкладки переключают \textbf{Operations},
\textbf{Live Tracking} и \textbf{Analytics}.

\subsection{Вид 1 --- Operations}
Обзор одним взглядом: \textbf{главная карточка скорости} (синяя на/выше оптимальной
скорости, красная ниже) со средней скоростью, числом кругов, временем заезда,
целевым кругом и дистанцией; сетка \textbf{Электрика}; сетка \textbf{Тепло \&
Эффективность}; и пульсирующая метка \textbf{LIVE}.

\subsection{Вид 2 --- Live Tracking \& карта}
Тёмная \textbf{OpenStreetMap} на весь экран (Leaflet), следящая за машиной.
\begin{itemize}
  \item \textbf{След, окрашенный по мощности} --- рисуется по сегментам и
    окрашивается по мощности двигателя: \textcolor{tgreen}{зелёный (0 Вт)} \arrow{}
    \textcolor{tamber}{янтарный (100 Вт)} \arrow{} \textcolor{tred}{красный
    (200 Вт+)}. Самая свежая часть непрозрачна и тускнеет со временем. (200 Вт ---
    максимальная мощность отслеживания ThaiGer; константа \code{POWER\_MAX}
    изменяема.)
  \item \textbf{Легенда мощности} (внизу слева) --- цветовая шкала и показание в
    ваттах в реальном времени.
  \item \textbf{Маркер машины} указывает направление движения.
  \item \textbf{Панель состояния GPS} (справа) --- качество фикса, точность, широта,
    долгота, высота, скорость GPS, курс, мощность двигателя, число кругов и счётчик
    \textbf{Filtered}.
\end{itemize}

\subsubsection{Очистка GPS-шума}
Каждый фикс очищается в четыре шага: (1) отбрасываются фиксы с точностью хуже 75 м;
(2) отклоняются «телепорты» --- фиксы с невозможной скоростью (более ~162 км/ч), с
защитой от ресинхронизации; (3) игнорируется дрожание менее 1,5 м на стоянке;
(4) лёгкое сглаживание. Результат --- чистый след без случайных скачков.

\subsubsection{Снимок за круг + сброс}
Каждый раз при увеличении счётчика кругов панель отрисовывает след завершённого
круга в PNG и скачивает его (\code{thaiger-lapN-<timestamp>.png}), затем очищает
след.
Зелёное уведомление подтверждает каждое сохранение.

\subsection{Вид 3 --- Analytics}
\textbf{Живая таблица телеметрии} со всеми 22 полями, а также график \textbf{Speed}
(км/ч во времени) и график \textbf{Efficiency} (КПД элемента и системы), каждый ---
по последним 60 точкам.

\section{Протокол телеметрии \& справочник полей}
Контроллер отправляет ASCII-кадры по BLE. Каждое поле обрамлено звёздочками с
однобуквенным ключом:

\begin{lstlisting}
*A28.3**B25.0**C3**D15:35**N53.4*
\end{lstlisting}
Парсер устойчив к потоку и всегда хранит последнее известное значение каждого поля.

\begin{tabularx}{\linewidth}{@{}clL@{}}
\toprule
\textbf{Ключ} & \textbf{Поле} & \textbf{Ед.}\\
\midrule
A & Скорость & км/ч\\
B & Средняя скорость & км/ч\\
C & Число кругов & ---\\
D & Общее время & мм:сс\\
E & Целевое время круга & мм:сс\\
F & Оптимальная скорость & км/ч\\
G & Напряжение элемента & В\\
H & Напряжение суперконденсатора & В\\
I & Напряжение двигателя & В\\
J & Ток элемента & А\\
K & Ток суперконденсатора & А\\
L & Ток двигателя & А\\
M & Собственное потребление & А\\
N & Температура элемента & \degC{}\\
O & Скважность воздушного насоса & \%\\
P & Подсказка по вождению & ---\\
Q & Разница напряжений ячеек & мВ\\
R & Энергия элемента (накопл.) & Вт·с\\
S & Энергия двигателя (накопл.) & Вт·с\\
T & КПД элемента & \%\\
U & КПД системы & \%\\
\bottomrule
\end{tabularx}

\medskip
Дистанция, мощность двигателя ($V \times A$) и энергия в Вт·ч вычисляются на
телефоне.

\section{Топики MQTT \& полезная нагрузка}
Все сообщения --- \textbf{QoS 0}. Пустые/NaN поля опускаются (нагрузка около
100--200 байт).

\noindent\textbf{\code{thaiger/\{car\}/telemetry}} --- публикуется с заданной
частотой (по умолчанию 5 Гц):

\begin{lstlisting}
{"ts":1736517201234,"A":28.3,"B":25.0,"C":3,"D":"15:35","G":28.3,"N":53.4,"dist":0.42}
\end{lstlisting}
\noindent\textbf{\code{thaiger/\{car\}/gps}} --- публикуется при наличии GPS-фикса:

\begin{lstlisting}
{"ts":1736517201300,"lat":13.756331,"lon":100.501762,"alt":12.0,"spd":53.4,"acc":3.5,"hdg":270.0}
\end{lstlisting}

\begin{tabularx}{\linewidth}{@{}lLl@{}}
\toprule
\textbf{Поле} & \textbf{Значение} & \textbf{Ед.}\\
\midrule
ts & Метка времени Unix & мс\\
lat / lon & Широта / долгота & град\\
alt & Высота (если есть) & м\\
spd & Скорость GPS (если есть) & км/ч\\
acc & Горизонтальная точность (если есть) & м\\
hdg & Курс (если есть) & град\\
\bottomrule
\end{tabularx}

\section{Профили автомобилей}

\begin{tabularx}{\linewidth}{@{}llllL@{}}
\toprule
\textbf{Профиль} & \textbf{id} & \textbf{Макс. мощн.} & \textbf{Предел темп.} & \textbf{Мин. цел. скорость}\\
\midrule
Thaiger 7 & \code{thaiger7} & 600 Вт & 70 \degC{} & 25 км/ч\\
Bengalo & \code{bengalo} & 900 Вт & 65 \degC{} & 25 км/ч\\
\bottomrule
\end{tabularx}

\medskip
Пороги можно переопределить для каждой машины в \textbf{Настройках}. След карты на
панели инженера использует отдельную шкалу мощности \textbf{200 Вт}
(\code{POWER\_MAX}).

\section{Устранение неполадок}

\begin{tabularx}{\linewidth}{@{}LL@{}}
\toprule
\textbf{Симптом} & \textbf{Причина / решение}\\
\midrule
Нет связи по Bluetooth & Сначала сопрягите модуль в настройках Bluetooth; проверьте питание; выдайте Bluetooth + Местоположение. Используйте \textbf{Pair New}.\\
«Connected», но нет данных & Не то устройство или контроллер молчит. Проверьте демо-режим, затем модуль.\\
Панель не зеленеет & Неверный WebSocket-URL (должен быть \code{wss://\ldots:8884/mqtt}, не порт \code{8883}), неверные данные или нет слушателя WebSocket.\\
Зелёная, но нет телеметрии & \textbf{Car} на панели не совпадает с id телефона, либо ретранслятор не включён. С обеих сторон одинаковый \code{thaiger7}/\code{bengalo}.\\
Пустая карта & Нет интернета для тайлов OSM или ещё нет GPS-фикса. Карта центрируется при первом фиксе.\\
След GPS скачет & Очистка должна это предотвращать; при очень грубых фиксах поднимите \code{GPS\_MAX\_ACC}. Счётчик \textbf{Filtered} показывает отброшенные.\\
GPS обновляется раз в несколько секунд & Обеспечьте открытое небо и разрешение \textbf{Местоположение}; в помещении --- грубые сетевые фиксы.\\
Нет CSV & Завершите заезд с записанными кадрами; иначе \emph{Export CSV} сообщает «No data to export».\\
\bottomrule
\end{tabularx}

\section{Тонкая настройка}
На телефоне (Настройки / Настройки MQTT): подсветка экрана, цвет скорости, частота
интерфейса, вибрация/звук, пороги по машине и все параметры брокера/ретранслятора.

\medskip
\noindent В \code{engineer\_dashboard.html} (константы в начале скрипта):

\begin{tabularx}{\linewidth}{@{}llL@{}}
\toprule
\textbf{Константа} & \textbf{По умолч.} & \textbf{Назначение}\\
\midrule
\code{POWER\_MAX} & 200 & Значение в ваттах для полного красного\\
\code{GPS\_MAX\_ACC} & 75 & Отбрасывать фиксы хуже этого числа метров\\
\code{GPS\_MAX\_SPEED} & 45 & м/с; быстрее --- считается телепортом\\
\code{GPS\_MIN\_MOVE} & 1.5 & м; меньшие перемещения --- дрожание на стоянке\\
\code{GPS\_SMOOTH} & 0.35 & Коэффициент сглаживания EMA (0 --- нет, 1 --- сырой)\\
\code{GPS\_MAX\_REJECTS} & 4 & Ресинхронизация после стольких отказов подряд\\
\code{TRAIL\_MAX} & 250 & Макс. число точек следа в памяти\\
\bottomrule
\end{tabularx}

\section{Глоссарий}
\begin{description}[style=nextline,leftmargin=0pt,labelwidth=0pt,labelsep=0pt,itemindent=0pt]
  \item[BLE] Bluetooth Low Energy, радиоканал от контроллера к телефону.
  \item[MQTT] Лёгкий протокол публикации/подписки; брокер ретранслирует сообщения.
  \item[Брокер] Сервер MQTT (например, HiveMQ Cloud), к которому подключаются обе стороны.
  \item[TLS] Шифрование соединения с брокером.
  \item[WebSocket] Транспорт, по которому браузер говорит на MQTT (\code{ws://} / \code{wss://}).
  \item[Топливный элемент (FC)] Водородный источник энергии на машине.
  \item[Суперконденсатор] Энергетический буфер для пиков нагрузки.
  \item[След (Trail)] Линия за машиной на карте, окрашенная по мощности двигателя.
\end{description}

\end{document}
